\subsection{多轴前缀与可核验语义：从单轴相位到多轴前缀的规范打包}\label{subsec:cdim-three-axis-address}
在“地址先于取值”的协议语义中，地址并非附加注释，而是读出\emph{可定义性}的前提：若地址对象未给定，则不存在可在对象层引用的评价函数，读出偏函数在该输入处不定义并以内禀空值 \(\mathsf{NULL}\) 表示（\S\ref{sec:recursive-addressing}；并见 \S\ref{subsec:xi-pw-type-safety-null} 的类型化口径）。
\par
为把该语义编译为可核验对象，令 \(A_{\mathrm{adm}}\) 为可接受地址集合，并对每个 \(a\in A_{\mathrm{adm}}\) 关联一个柱事件 \(C_a\)（cylinder event）；细化关系 \(b\lhd a\) 对应柱事件包含 \(C_b\subseteq C_a\)，从而 \((A_{\mathrm{adm}},\lhd)\) 形成前缀偏序生成的站点。
在该站点上，把“协议允许且在类型意义下为单值可核验的局部证书对象”组织为（预）层 \(F\)：对每个地址 \(a\)，令 \(F(a)\) 为在 \(C_a\) 上可接受的局部证书集合；若 \(b\lhd a\) 则存在限制映射 \(F(a)\to F(b)\) 表示证书细化/投影（参见定义 \ref{def:xi-cert-refinement-equivalence} 的可核验等价与共同细化口径）。

\begin{theorem}[有限前缀不导出唯一识别：可计算概念在同阶描述长度下的前缀不可判别性]\label{thm:cdim-computable-prefix-nonidentifiability}
\leanverified{paper\_cdim\_computable\_prefix\_nonidentifiability}
固定一台前缀自由通用机 \(U\)。
对可计算无穷二元序列 \(x\in\{0,1\}^{\NN}\)，定义其生成复杂度为
\[
K_{\mathrm{gen}}(x):=\min\Bigl\{|p|:\ \forall n\ge 1,\ U(p,n)=x\upharpoonright n\Bigr\},
\]
其中 \(U(p,n)\) 表示程序 \(p\) 在输入 \(n\) 时输出 \(x\) 的前 \(n\) 位。
则对任意 \(n\ge 1\)，存在另一可计算序列 \(y\neq x\) 满足
\[
y\upharpoonright n=x\upharpoonright n,
\qquad
K_{\mathrm{gen}}(y)\ \le\ K_{\mathrm{gen}}(x)+O(\log n).
\]
因此，在可计算概念类内，任何有限前缀一致性都不能推出唯一识别；解除歧义需要额外的对数级索引/证书，而非继续读取同一前缀轴本身。
\end{theorem}
\begin{proof}
取实现 \(K_{\mathrm{gen}}(x)\) 的最短程序 \(p_x\)。
构造程序 \(p_y\)：在输入 \(m\) 时，先调用 \(p_x\) 计算 \(x\upharpoonright \max\{m,n+1\}\)；
随后输出前 \(n\) 位保持不变，将第 \(n+1\) 位取反，并把其后位（若 \(m>n+1\)）置为 \(0\)。
该过程给出一条可计算序列 \(y\)，显然 \(y\upharpoonright n=x\upharpoonright n\) 且 \(y\neq x\)。
程序 \(p_y\) 仅需在 \(p_x\) 的基础上额外编码整数 \(n\) 与固定改写逻辑，故
\(
|p_y|\le |p_x|+O(\log n)
\)，即得所示不等式。证毕。
\end{proof}

\begin{theorem}[有限前缀不可判别性的任意有限多重化]\label{thm:cdim-computable-prefix-nonidentifiability-multiplicity}
在定理 \ref{thm:cdim-computable-prefix-nonidentifiability} 的设定下，对任意可计算无穷二元序列 \(x\in\{0,1\}^{\NN}\)、任意 \(n\ge 1\) 与任意整数 \(t\ge 1\)，存在两两不同的可计算序列
\[
y^{(1)},\dots,y^{(t)}
\]
满足
\[
y^{(j)}\upharpoonright n=x\upharpoonright n
\qquad
(1\le j\le t),
\]
且
\[
K_{\mathrm{gen}}(y^{(j)})
\le
K_{\mathrm{gen}}(x)+O(\log n+\log t)
\qquad
(1\le j\le t).
\]
亦即，共享长度 \(n\) 前缀的语义不可判别类，可以在仅增加 \(O(\log n+\log t)\) 描述开销的条件下扩张到任意预先指定的有限规模 \(t\)。
\leanverified{paper\_cdim\_computable\_prefix\_nonidentifiability\_multiplicity}
\end{theorem}
\begin{proof}
固定实现 \(K_{\mathrm{gen}}(x)\) 的最短程序 \(p_x\)。对每个 \(1\le j\le t\)，构造程序 \(p_j\) 如下：在输入 \(m\) 时，先调用 \(p_x\) 计算 \(x\upharpoonright \max\{m,n+j\}\)；随后定义输出序列 \(y^{(j)}\) 为
\[
y^{(j)}(r)=
\begin{cases}
x(r),& r<n+j,\\
1-x(n+j),& r=n+j,\\
0,& r>n+j.
\end{cases}
\]
于是每个 \(y^{(j)}\) 都是可计算序列，并且其前 \(n\) 位与 \(x\) 完全相同。若 \(i\neq j\)，则 \(y^{(i)}\) 与 \(y^{(j)}\) 在位置 \(n+\min\{i,j\}\) 处已经分离，故两两不同。
\par
为了生成 \(y^{(j)}\)，程序 \(p_j\) 只需在 \(p_x\) 的基础上额外编码整数 \(n\)、整数 \(j\) 以及一段固定长度的修改规则。因此
\[
|p_j|
\le
|p_x|+O(\log n+\log j)
\le
|p_x|+O(\log n+\log t).
\]
取最短描述长度即得
\[
K_{\mathrm{gen}}(y^{(j)})
\le
K_{\mathrm{gen}}(x)+O(\log n+\log t),
\]
结论成立。证毕。
\end{proof}
这说明前缀预算不足所导致的语义歧义并非只表现为“存在一个竞争者”，而是能够以对数额外开销支持任意给定规模的有限歧义簇。

\begin{proposition}[\texorpdfstring{$\mathsf{NULL}$}{NULL} 的对象论判据：实际全局与兼容局部]\label{prop:cdim-nonnull-section-criterion}
在上述站点/（预）层表述下，记 \(F^\sharp\) 为 \(F\) 的 sheafification。则：
\begin{enumerate}
  \item 地址 \(a\) 的读出在类型层非 \(\mathsf{NULL}\) 当且仅当
  \[
  \boxed{\ F(a)\neq\varnothing\ }.
  \]
  \item 有限柱覆盖上的兼容局部对象族按共同 refinement 取商后，恰好给出 \(F^\sharp(a)\) 的元素；因此
  \[
  F^\sharp(a)\neq\varnothing
  \]
  当且仅当地址 \(a\) 处存在兼容局部截面族。
  \item 只有在 \(F\) 于 \(a\) 处已经满足层条件时，\(F^\sharp(a)\neq\varnothing\) 才等价于 \(F(a)\neq\varnothing\)。因此“兼容局部存在”一般不能直接读作“原始对象层已有全局读出”。
\end{enumerate}
\end{proposition}
\begin{proof}
第一条是 \(\mathsf{NULL}\) 的定义性译文：在地址 \(a\) 上，“读出可定义且可核验”被组织为存在某个可接受局部证书对象。
\par
第二条只是 sheafification 的标准匹配族描述：相容性以“在重叠上存在共同细化”表达（定义 \ref{def:xi-cert-refinement-equivalence}），并在共同 refinement 关系下组成 \(F^\sharp(a)\) 的元素。
\par
第三条则是层化条件在点 \(a\) 处的直接推论：若 \(F\to F^\sharp\) 在 \(a\) 处为双射，则兼容局部匹配族与原始对象层全局证书严格一致；反之二者必须分开记录。证毕。
\end{proof}

\paragraph{三轴地址的规范打包}
在第 \ref{subsec:xi-pw-type-safety-null} 节的统一口径下，一个可核对地址被规范打包为多轴前缀
\[
a=(a_{\mathrm{vis}},a_{\mathrm{res}},a_{\mathrm{mode}})
\]
（定义 \ref{def:multi_axis_address}）。
其中 \(a_{\mathrm{vis}}\) 对齐 unitary slice 的单参相位接口（相位前缀/柱事件），\(a_{\mathrm{res}}\) 承载为避免同像碰撞所需的 residue/逆极限轴外置（prime-register），而 \(a_{\mathrm{mode}}\) 记录协议所允许的正交模式轴（典型包含径向寄存器与最小实现维数轴）；该三轴打包把“地址对象”从单轴相位提升为可编译的最小接口直和，并为后续的预算—失败类型分离提供规范坐标。
在“统一最小性 + 协议闭包（无未记录轴回流）”假设下，模式轴的完备性可由 Hankel 秩一致有界刻画（定理 \ref{thm:xi-completeness-iff-hankel-bounded}），其可实现的离线核验接口可取为 ACC-Rank 稳定性证书（注 \ref{rem:xi-acc-rank}）。

\subsection{\texorpdfstring{$\mathsf{NULL}$}{NULL} 的三分解与失败证书：语义空值、协议空值、碰撞空值}\label{subsec:cdim-null-three-way}
在站点/（预）层表述下，\(\mathsf{NULL}\) 的来源需区分为“局部不存在”“兼容局部仅停留在 sheafification”“以及单值唯一化失败”三层；并且该失败在体系内可被分解为互不混淆的三类离线可核验见证（定理 \ref{thm:xi-null-complete-trichotomy-offline}）。
与三轴地址语义对齐，可将其表述为：
\begin{enumerate}
  \item \textbf{语义空值：}地址未给定或 \(a\notin A_{\mathrm{adm}}\)。此时 \(C_a\) 不存在（指称缺失），故 \(F(a)\) 无从定义或按约定为空，从而读出为 \(\mathsf{NULL}\)。
  \item \textbf{协议空值：}地址 \(a\in A_{\mathrm{adm}}\) 但协议/类型系统拒绝在 \(C_a\) 上建立读出（典型为 unitary slice 锁定与 \(\mathbf{PW}\) 闭包下的“无未记录连续轴”原则，定理 \ref{thm:xi-pw-no-continuous-hair}），或因模式轴/最小实现维数门槛被违反而触发维数型 \(\mathsf{NULL}\)（定理 \ref{thm:xi-completeness-iff-hankel-bounded}；注 \ref{rem:xi-acc-rank}），或虽有兼容局部对象 \(F^\sharp(a)\neq\varnothing\) 却因二阶一致性失败而无法下降为原始对象层全局读出；其失败可由 \(H^2\) 粘合障碍类 \([\alpha]\neq 0\) 给出可核验见证（命题 \ref{prop:null-as-h2-obstruction}）。
  \item \textbf{碰撞空值：}地址 \(a\in A_{\mathrm{adm}}\) 且通过检查器，但底层读出存在多对一压缩而未外置足够残差轴使扩展读出单射化，或等价地存在不同全局候选在同一有限覆盖上给出不可区分的局部限制；其最小侧信息预算由地址—缺陷定律（定理 \ref{thm:address-defect-law}）与分位配额阈值（推论 \ref{cor:pom-prime-register-quantile-budget}）给出；端点侧则由二次分辨率屏障（定理 \ref{thm:xi-offcritical-endpoint-resolution-lower-bound}）强制。
\end{enumerate}
上述三类失败均可输出有限、离线可核验的否定证书（地址缺失/类型拒绝/预算不足或障碍类见证），从而把“不可核验”转换为可证伪接口。

\subsection{正交阈值的不可替代性：预算不能互相抵偿的结构定理}\label{subsec:cdim-orthogonal-threshold-nonsubstitutable}
在 unitary slice 可见域约束与 \(\mathbf{PW}\) 闭包口径下，失败来源在体系内分解为互不替代的正交见证（定理 \ref{thm:xi-null-orthogonal-trichotomy-resource-nonsubstitutable}）。
该不可替代性在三轴地址语义下可解释为：不同失败类型由不同记录轴支撑，其预算门槛位于互不重叠的正交通道上，因而不存在用一类预算覆盖另一类缺口的合法编译路径。

\begin{theorem}[预算—失败类型分离（正交阈值不可替代性）]\label{thm:cdim-budget-failure-separation}
在三轴地址 \(a=(a_{\mathrm{vis}},a_{\mathrm{res}},a_{\mathrm{mode}})\) 与离线编译—核验流程（算法 \ref{alg:cdim-three-axis-address-compile}）下，\(\mathsf{NULL}\) 的三类来源（\S\ref{subsec:cdim-null-three-way}）具有如下分离性质：
\begin{enumerate}
  \item 若失败为语义空值（\(a\notin A_{\mathrm{adm}}\)），则任意预算提升均无意义；该失败仅可通过改变指称/地址对象消除。
  \item 若失败为协议空值（类型拒绝或 \(H^2\) 粘合障碍 \([\alpha]\neq 0\)），则任何 residue/端点/绕数等预算提升均不能消去该失败见证（命题 \ref{prop:null-as-h2-obstruction}）。
  \item 若失败为碰撞空值（单射化预算不足），则对任意给定置信口径，其修复\emph{必须}发生在相应支撑轴上：prime-register 预算须越过分位门槛（推论 \ref{cor:pom-prime-register-quantile-budget}）且端点分辨率须满足二次门槛（定理 \ref{thm:xi-offcritical-endpoint-resolution-lower-bound}）；增加其它正交轴预算不能替代该跨越。
\end{enumerate}
因此，在该体系内“预算”与“失败类型”严格正交：每一类失败都有独立的阈值/不变量证书，且预算提升不能跨轴抵偿。
\end{theorem}
\leanverified{paper\_cdim\_budget\_failure\_separation}
\begin{proof}
语义空值、协议空值与碰撞空值的完备性与可离线核验性由定理 \ref{thm:xi-null-complete-trichotomy-offline} 给出；其中协议空值的“不可用预算消去”由 \(H^2\) 障碍见证（命题 \ref{prop:null-as-h2-obstruction}）给出，碰撞空值的“必须跨越相应阈值”由 prime-register 分位门槛（推论 \ref{cor:pom-prime-register-quantile-budget}）与端点二次门槛（定理 \ref{thm:xi-offcritical-endpoint-resolution-lower-bound}）联合强制；三类见证的正交与不可替代性由定理 \ref{thm:xi-null-orthogonal-trichotomy-resource-nonsubstitutable} 封口。证毕。
\end{proof}
例如：
\begin{itemize}
  \item prime-register 的模积预算只约束碰撞方向的可逆化（定理 \ref{thm:address-defect-law} 与推论 \ref{cor:pom-prime-register-quantile-budget}），不能抵消端点二次分辨率门槛（定理 \ref{thm:xi-offcritical-endpoint-resolution-lower-bound}）；
  \item 任意预算提升都不能消去粘合障碍类 \([\alpha]\neq 0\) 的同调见证（命题 \ref{prop:null-as-h2-obstruction}），其失败属于结构性不可全球化而非资源不足；
  \item 若试图引入与相位解耦的额外连续坐标以绕过 \(\mathsf{NULL}\)，则在“无未记录连续轴”原则下只能被类型拒绝或被迫因子化到唯一径向寄存器 \(\varrho=\log|u|\)（推论 \ref{cor:xi-unique-continuous-transverse-register} 与定理 \ref{thm:xi-pw-no-continuous-hair}）。
\end{itemize}
因此，“信息守恒/不泄漏”的可审计含义被落实为：失败来源在正交分解下各自闭合，并分别携带独立阈值或不变量证书；预算提升只能在其所属轴上发生有效补偿，不能跨轴抵偿。

\paragraph{前缀深度与单射化预算}
相位前缀过滤在端点紧化下所诱导的二次分辨率门槛见 \S\ref{subsec:cdim-phase-prefix-quadratic-barrier}。
碰撞方向的单射化预算阈值由 prime-register 分位配额律给出（推论 \ref{cor:pom-prime-register-quantile-budget}；并见 \S\ref{subsec:xi-pw-type-safety-null}）。
因此，在多轴前缀口径下，\(a_{\mathrm{vis}}\) 的位深不足只能通过 \(a_{\mathrm{res}}\) 与其它正交证书轴外置修复；拒绝外置将以碰撞型 \(\mathsf{NULL}\) 表达。

\subsection{可直接编译的判定流程：从地址 \texorpdfstring{$a$}{a} 到证书对象或 \texorpdfstring{$\mathsf{NULL}$}{NULL} 见证}\label{subsec:cdim-address-compiler}
在上述结构下，可将“可编译、可核验、可证伪”的证书协议规定为一个完全离线可核验的编译—检查流程；其输入为三轴地址 \(a=(a_{\mathrm{vis}},a_{\mathrm{res}},a_{\mathrm{mode}})\)，输出为一个证书对象 \(s\in F(a)\) 或一份 \(\mathsf{NULL}\) 失败见证。

\begin{algorithm}[三轴地址的编译—核验流程（离线可复核）]\label{alg:cdim-three-axis-address-compile}
在站点 \(A_{\mathrm{adm}}\) 与（预）层 \(F\) 的设定下，定义如下流程：
\begin{enumerate}
  \item \textbf{地址合法性检查：}验证 \(a\in A_{\mathrm{adm}}\)。若失败则输出语义空值证书并返回 \(\mathsf{NULL}\)。
  \item \textbf{类型/可见域检查：}验证协议允许在 \(C_a\) 上建立读出（unitary slice 锁定与模式轴许可；定理 \ref{thm:xi-pw-no-continuous-hair}）。若失败则输出协议空值证书并返回 \(\mathsf{NULL}\)。
  \item \textbf{单射化预算与障碍检查：}对 \(a_{\mathrm{mode}}\) 核验模式轴完备性（Hankel 秩一致有界或等价的 ACC-Rank 稳定性，定理 \ref{thm:xi-completeness-iff-hankel-bounded}；注 \ref{rem:xi-acc-rank}）；对 \(a_{\mathrm{res}}\) 核验 prime-register 的模积阈值（推论 \ref{cor:pom-prime-register-quantile-budget} 或定理 \ref{thm:address-defect-law} 的硬下界）；对端点精度预算核验二次分辨率门槛（定理 \ref{thm:xi-offcritical-endpoint-resolution-lower-bound}）；对绕数/同调障碍核验相应不变量是否为零（定理 \ref{thm:dephys-dirichlet-winding-minimal}，命题 \ref{prop:null-as-h2-obstruction}）。若任一不足则输出对应失败见证并返回 \(\mathsf{NULL}\)。
  \item \textbf{对象构造与 descent 检查：}在有限柱覆盖上构造局部证书对象族并核验交叠相容性；若相容则先得到 sheafified 兼容类 \(s^\sharp\in F^\sharp(a)\)。只有在附加的 sheaf-at-\(a\) 或等价 descent 证书通过时，才把 \(s^\sharp\) 下降为唯一全局对象 \(s\in F(a)\) 并返回；否则输出协议空值见证（命题 \ref{prop:cdim-nonnull-section-criterion}）。
\end{enumerate}
\end{algorithm}
该流程与最小记录轴定理一致：任何为消除 \(\mathsf{NULL}\) 而引入的新轴，必须在 \(\mathrm{Ext}(\pi)\) 中被识别为规范初对象 \(r_0\) 的因子化；而在连续轴层面，除径向寄存器 \(\varrho(u)=\log|u|\) 的等价类外不存在可接受扩维通道（定理 \ref{thm:cdim-minimal-record-axis}）。

\subsection{边界层盲区闭环：压缩、盲区与外置前缀/可逆层析}\label{subsec:cdim-boundary-blindspot-loop}
在有限证书口径下，Jensen 缺陷 \(D_{\zeta}(\varrho)\) 作为 \(\log\varrho\) 的函数在 \(\varrho=|w_\rho|\) 处出现折角原子，其斜率跳变等于相应零点的重数（推论 \ref{cor:app-jensen-defect-kink-structure}）。
因此，任何仅使用有限半径采样族的核验策略都存在结构性盲区：若采样半径上界未越过 \(|w_\rho|\)，则该折角阈值在采样上严格不可见，无法排除相应零点（推论 \ref{cor:app-jensen-finite-sampling-blindspot}）。
\par
在端点二次压缩制度下，盲区与位预算门槛同阶出现：固定可见接口若拒绝外置可寻址前缀，则离切片偏移在高高度区间被迫通过端点二次径向通道进入并触发不可约的端点分辨率/位复杂度下界（定理 \ref{thm:xi-offcritical-endpoint-resolution-lower-bound}；并见注 \ref{rem:xi-comoving-horizon-localization-address-null} 对制度成本的说明）。
若允许外置可寻址前缀并取共动局部化，则二次红移可被解除为线性显影，端点层位预算降为与高度窗无关的常数门槛（推论 \ref{cor:xi-comoving-prefix-bit-budget-gain}）。
在本章的圆盘完成化口径下，幂覆盖/Adams--Norm 折叠门提供了与共动前缀同型的“解压”机制：它把端点边界层中的点状缺陷搬运到预设紧子圆盘并在固定半径处以 Jensen 缺陷证书显影（推论 \ref{cor:xi-adams-norm-fixed-radius-jensen-certificate}；并见 \S\ref{subsec:cdim-adams-norm-folding-jensen}）。
\par
上述结构形成一个实现无关的闭环：
\begin{itemize}
  \item 固定图表的端点压缩把离切片信息推入边界层，从而产生有限证书的不可判别盲区；
  \item 解除盲区的两条可审计路径分别是：外置可寻址前缀以共动局部化（消除二次红移），或进入可逆层析通道（把二维参数编码为一维指数谱并可反演，定理 \ref{thm:cdim-comoving-horizon-scan-fourier-inversion}）；
  \item 因而“地址先于取值”的制度后果可被定量化为：是否外置前缀决定了位预算主项是否包含 \(2\log_2T\)，并决定了边界层盲区是否可被穿透（参见 \S\ref{subsec:discussion-comoving-tomography-prefix-budget-blindspot} 的推广表述）。
\end{itemize}

\subsection{共动缺陷的降红移机制：从二次端点压缩到一阶可见性}\label{subsec:cdim-comoving-defect-deredshift}
\subsubsection{共动图表与共动 Jensen 缺陷}
固定图表制度下，离切片点状缺陷在 Cayley--Busemann 机制中被压入端点边界层并携带二次红移；因此仅以固定半径的 Jensen 缺陷作证书时，排除高度窗上的偏移必须支付与 \(\gamma^2\) 同阶增长的端点分辨率预算（定理 \ref{thm:xi-offcritical-endpoint-resolution-lower-bound}）。
与之相对，移位 Cayley 图表所诱导的平移族完成化 \(F_{\zeta,x_0}\) 及其 Jensen 缺陷 \(D_{\zeta,x_0}\)（式 \eqref{eq:xi-comoving-shifted-Fzeta}--\eqref{eq:xi-comoving-shifted-jensen-defect}）允许把“高度红移”转写为图表选择：在共动选择 \(x_0=\gamma\) 下，离切片零点 \(\rho=\beta+\mathrm{i}\gamma\)（\(\beta<\tfrac12\)）的圆盘像退化为
\[
w_{\rho,\gamma}=\mathcal{C}(\mathrm{i}\delta)=\frac{1-\delta}{1+\delta}\in(0,1),
\qquad
\delta:=\tfrac12-\beta>0
\]
（式 \eqref{eq:xi-comoving-pure-delta}），从而 \(|w_{\rho,\gamma}|\) 仅依赖偏移 \(\delta\) 而与高度无关。

\subsubsection{定量偏移上界：共动缺陷上界推出偏移的显式控制}
\begin{proposition}[共动缺陷界导出无高度偏移界]\label{prop:cdim-comoving-defect-implies-delta-bound}
对任意 \(x_0\in\RR\) 令 \(D_{\zeta,x_0}\) 为平移族完成化的 Jensen 缺陷（式 \eqref{eq:xi-comoving-shifted-jensen-defect}）。
若存在 \(0<\varrho<1\) 与 \(\varepsilon\ge 0\) 使
\[
D_{\zeta,x_0}(\varrho)\le \varepsilon,
\]
则对任意离切片零点 \(\rho=\beta+\mathrm{i}\gamma\) 满足 \(\beta<\tfrac12\) 且 \(\gamma=x_0\) 时，其偏移 \(\delta=\tfrac12-\beta\) 满足显式上界
\[
\boxed{\;
\delta\ \le\ \frac{1-\varrho e^{-\varepsilon}}{1+\varrho e^{-\varepsilon}}\; }.
\]
\leanverified{paper\_cdim\_comoving\_defect\_implies\_delta\_bound\_seeds}
\end{proposition}
\begin{proof}
该命题即定理 \ref{thm:xi-comoving-defect-implies-delta-bound} 在本章符号下的直接表述；其证明以单零点 Jensen 下界（命题 \ref{prop:app-jensen-single-zero-lower-bound}）与共动像 \(|w_{\rho,\gamma}|=\frac{1-\delta}{1+\delta}\) 的代数消元为核心。证毕。
\end{proof}
该结论的制度含义在于：在保持 unitary slice 的单参可见轴不变的前提下，外置共动前缀可将“端点二次压缩导致的不可见性”转写为半径阈值上的确定性比较，从而把偏移排除从 \(\gamma^2\)-量纲降为 \(\delta\)-量纲的直接估计（并与共动前缀的位预算收益律相容，推论 \ref{cor:xi-comoving-prefix-bit-budget-gain}）。

\subsubsection{一致小化判别：以 \texorpdfstring{$\sup_{x_0}$}{sup\_{x0}} 消除图表选择}
共动选择 \(x_0=\gamma\) 的制度含义在于：对任意潜在离切片零点，都存在一条图表使其偏移以一阶尺度进入半径阈值（式 \eqref{eq:xi-comoving-pure-delta}）。
将该“可随零点移动的图表选择”提升为实现无关的统一证书语法，可得到 RH 的一致小化判别：通过对 \(x_0\) 取上确界，将图表选择消去为一个完全离线可核验的量词结构。
\begin{theorem}[RH 的共动 Jensen 缺陷一致小化判别]\label{thm:cdim-rh-iff-uniform-comoving-defect}
下列命题等价：
\begin{enumerate}
  \item[(i)] RH 成立；
  \item[(ii)] 存在 \(\varrho_n\uparrow 1\)、\(\varepsilon_n\downarrow 0\)，使得对所有 \(n\) 有一致界
  \[
  \sup_{x_0\in\RR}D_{\zeta,x_0}(\varrho_n)\le \varepsilon_n.
  \]
\end{enumerate}
\end{theorem}
\begin{proof}
该定理即定理 \ref{thm:xi-rh-iff-uniform-comoving-defect} 在本章符号下的直接译文。证毕。
\leanverified{paper\_cdim\_rh\_iff\_uniform\_comoving\_defect}
\end{proof}

\paragraph{共动全局积分缺陷与阈值尖点（引用）}
共动 Jensen 缺陷在平移基点方向上的全局积分分解、阈值尖点律与黄金分辨率链渐近，统一见段落 \ref{par:xi-comoving-integrated-jensen-defect-sumrule} 及定理 \ref{thm:xi-integrated-defect-cusp-reconstruction}。

\subsection{共动视界扫描的傅里叶完备性：有限缺陷口径下一维剖面的可逆层析}\label{subsec:cdim-comoving-horizon-scan-fourier-completeness}
设存在有限离临界线零点多重集 \(\{\rho_j\}_{j=1}^{M}\)（按重数计），其中
\(\rho_j=\beta_j+\mathrm{i}\gamma_j\) 且 \(\beta_j<\tfrac12\)。
令 \(\delta_j:=\tfrac12-\beta_j>0\)，并记重数 \(m_j\in\ZZ_{\ge 1}\)。
对任意共动参数 \(x_0\in\RR\)，定义移位图表下的圆盘像
\[
w_{\rho_j,x_0}:=\mathcal{C}\!\bigl((\gamma_j-x_0)+\mathrm{i}\delta_j\bigr)\in\mathbb{D},
\qquad
h_{x_0}(\rho_j):=1-|w_{\rho_j,x_0}|^2,
\]
并定义聚合剖面函数
\[
H(x_0):=\sum_{j=1}^{M}m_j\,h_{x_0}(\rho_j).
\]
由附录定理 \ref{thm:app-comoving-horizon-localization} 的闭式表达，\(h_{x_0}(\rho)\) 是以 \((\gamma,\delta)\) 为中心/宽度参数的洛伦兹窗：
\[
\boxed{\;
h_{x_0}(\rho_j)=\frac{4\delta_j}{(\gamma_j-x_0)^2+(1+\delta_j)^2}\;}
\]
从而 \(H\) 为有限 Cauchy/Lorentz 核混合。

\begin{theorem}[扫描剖面的 Fourier--Laplace 指纹：有限指数混合与可辨识性]\label{thm:cdim-comoving-horizon-scan-fourier-inversion}
在上述设定下：
\begin{enumerate}
  \item \(H\in L^1(\RR)\cap C^\omega(\RR)\)；
  \item 其 Fourier 变换具有显式谱分解：对任意 \(\xi\in\RR\)，令
  \[
  \widehat H(\xi):=\int_{\RR}e^{-\mathrm{i}\xi x_0}H(x_0)\,\dd x_0,
  \]
  则
  \[
  \boxed{\;
  \widehat H(\xi)
  =4\pi\sum_{j=1}^{M}m_j\,\frac{\delta_j}{1+\delta_j}\,e^{-(1+\delta_j)|\xi|}\,e^{-\mathrm{i}\gamma_j\xi}\; } ;
  \]
  \item 若两组有限离线零点多重集诱导同一剖面 \(H\)（只需在任意非平凡开区间上一致），则两组多重集 \(\{(\gamma_j,\delta_j,m_j)\}\) 按重数一致。
\end{enumerate}
\end{theorem}
\begin{proof}
该定理即定理 \ref{thm:xi-comoving-horizon-scan-fourier-inversion} 在本章符号下的直接译文。证毕。
\end{proof}
该结论把二维参数 \((\gamma,\delta)\) 压缩为一维剖面 \(x_0\mapsto H(x_0)\) 的指数谱：\(\gamma\) 以平移协变的相位载体 \(e^{-\mathrm{i}\gamma\xi}\) 出现，而 \(\delta\) 以 \(e^{-(1+\delta)|\xi|}\) 的 Laplace 衰减出现；从而在有限缺陷口径下得到完全可逆的层析编码。

\begin{theorem}[最近缺陷层的 Fourier--Laplace 首层提取]\label{thm:cdim-comoving-horizon-scan-first-layer-extraction}
在定理 \ref{thm:cdim-comoving-horizon-scan-fourier-inversion} 的设定下，把不同深度值按升序排列为
\[
0<\delta^{(1)}<\delta^{(2)}<\cdots<\delta^{(r)},
\]
并定义首层指标集
\[
I_1:=\{j:\delta_j=\delta^{(1)}\}.
\]
若 \(r=1\)，约定 \(\delta^{(2)}:=+\infty\)，从而相应误差项视为 \(0\)。
则当 \(|\xi|\to\infty\) 时成立精确首层分离
\[
\boxed{\;
e^{(1+\delta^{(1)})|\xi|}\widehat H(\xi)
=
G_1(\xi)+
O\!\bigl(e^{-(\delta^{(2)}-\delta^{(1)})|\xi|}\bigr)\;}
\]
其中
\[
G_1(\xi):=
4\pi\sum_{j\in I_1}m_j\frac{\delta^{(1)}}{1+\delta^{(1)}}e^{-\mathrm{i}\gamma_j\xi}
\]
为有限三角多项式；若 \(r=1\)，则上式中的误差项按约定为 \(0\)。
特别地，最近临界层 \(\delta^{(1)}\) 及其全部横向位置 \(\{\gamma_j\}_{j\in I_1}\) 与重数 \(\{m_j\}_{j\in I_1}\) 可由高频尾部唯一恢复。
\end{theorem}
\begin{proof}
由定理 \ref{thm:cdim-comoving-horizon-scan-fourier-inversion} 的显式公式，按不同深度分组得
\[
\widehat H(\xi)
=
4\pi\sum_{\nu=1}^{r}
e^{-(1+\delta^{(\nu)})|\xi|}
\sum_{j:\delta_j=\delta^{(\nu)}}
m_j\frac{\delta^{(\nu)}}{1+\delta^{(\nu)}}e^{-\mathrm{i}\gamma_j\xi}.
\]
提出最小深度层 \(\delta^{(1)}\)，得到
\[
e^{(1+\delta^{(1)})|\xi|}\widehat H(\xi)
=
G_1(\xi)+
\sum_{\nu=2}^{r}
e^{-(\delta^{(\nu)}-\delta^{(1)})|\xi|}
G_\nu(\xi),
\]
其中各 \(G_\nu\) 仍为有限三角多项式，因而一致有界。若 \(r\ge 2\)，余项由最小正差
\(\delta^{(2)}-\delta^{(1)}\) 控制，遂得所示估计；若 \(r=1\)，则根本不存在余项。
\par
再证唯一恢复性。把 \(G_1\) 中具有相同横向位置的项合并，可写成
\[
G_1(\xi)=\sum_{\ell=1}^{r_1} c_\ell e^{-\mathrm{i}\eta_\ell\xi},
\qquad
c_\ell>0,
\]
其中 \(\eta_\ell\) 为首层不同 \(\gamma_j\) 的集合。有限频率指数函数族在 \(\RR\) 上线性无关，故该表示唯一；而
\[
c_\ell=
4\pi\frac{\delta^{(1)}}{1+\delta^{(1)}}
\sum_{j\in I_1:\,\gamma_j=\eta_\ell}m_j
\]
又唯一确定对应的总重数。于是首层深度、横向位置与重数全部由高频尾部唯一恢复。证毕。
\end{proof}
这说明共动视界扫描的高频尾部并非只给出“某种可辨识性”，而是带有天然的层级结构：最靠近临界线的一层总是先被读出，其后的各层只以更快衰减项出现。

\subsection{二维 Fourier--Laplace 全息层析：边界指纹唯一恢复缺陷测度}\label{subsec:cdim-fourier-laplace-holographic-tomography}
将有限多重集推广为 \(\RR\times(0,\infty)\) 上有限正测度 \(\nu\)。
定义边界密度
\[
H_\nu(x):=\int_{\RR\times(0,\infty)}\frac{4\delta}{(x-\gamma)^2+(1+\delta)^2}\,\dd\nu(\gamma,\delta),
\qquad
H_{\nu,t}:=P_t*H_\nu,
\]
其中 \(P_t(x)=\frac{1}{\pi}\frac{t}{x^2+t^2}\) 为 Poisson 核。
记 \(\widehat H_{\nu,t}(\xi)\) 为 Fourier 变换，令 \(s:=|\xi|\) 并定义二参数指纹
\[
F(\xi,s):=e^{(t+1)s}\widehat H_{\nu,t}(\xi).
\]

\begin{proposition}[二维指纹的张量化分解与单射性]\label{prop:cdim-defect-measure-fourier-laplace-holographic}
在上述设定下成立严格分解
\[
\boxed{\;
F(\xi,s)=4\pi\int_{\RR\times(0,\infty)}\frac{\delta}{1+\delta}\,e^{-s\delta}\,e^{-\mathrm{i}\gamma\xi}\,\dd\nu(\gamma,\delta)\; } ,
\]
从而 \(F\) 是 \(\gamma\) 的 Fourier 变换与 \(\delta\) 的 Laplace 变换之张量化组合；特别地映射 \(\nu\mapsto F\) 在有限测度类上为单射，因而缺陷测度 \(\nu\) 可由边界可见指纹 \(F\) 的取值唯一恢复。
\end{proposition}
\begin{proof}
该命题即结论 \ref{con:xi-terminal-defect-measure-fourier-laplace-holographic} 的直接表述。证毕。
\end{proof}

\begin{theorem}[有限缺陷测度的递归层剥离唯一性]\label{thm:cdim-defect-measure-recursive-layer-peeling}
在命题 \ref{prop:cdim-defect-measure-fourier-laplace-holographic} 的设定下，若
\[
\nu=\sum_{j=1}^{\kappa} m_j\,\delta_{(\gamma_j,\delta_j)}
\qquad
(m_j\in\ZZ_{\ge 1})
\]
为有限原子缺陷测度，并把不同深度值按升序排列为
\[
0<\delta^{(1)}<\delta^{(2)}<\cdots<\delta^{(r)},
\]
并约定 \(\delta^{(r+1)}:=+\infty\)。
则 \(\nu\) 可由 \(\widehat H_\nu(\xi)\) 通过有限步递归剥离唯一恢复。更精确地，若前 \(s-1\) 层已恢复并定义
\[
\nu^{(<s)}:=
\sum_{\nu'=1}^{s-1}\ \sum_{j:\delta_j=\delta^{(\nu')}}m_j\,\delta_{(\gamma_j,\delta_j)},
\]
以及残差 Fourier 变换
\[
\widehat H^{(s)}(\xi):=\widehat H_\nu(\xi)-\widehat H_{\nu^{(<s)}}(\xi),
\]
则对 \(1\le s\le r\) 有
\[
e^{(1+\delta^{(s)})|\xi|}\widehat H^{(s)}(\xi)
=
G_s(\xi)+
O\!\bigl(e^{-(\delta^{(s+1)}-\delta^{(s)})|\xi|}\bigr),
\]
其中
\[
G_s(\xi)=
4\pi\sum_{j:\delta_j=\delta^{(s)}}
m_j\frac{\delta^{(s)}}{1+\delta^{(s)}}e^{-\mathrm{i}\gamma_j\xi},
\]
特别地，当 \(s=r\) 时上式中的误差项自动视为 \(0\)。因此经过 \(r\) 步递归后，整个有限测度 \(\nu\) 被唯一确定。
\end{theorem}
\begin{proof}
由 \(H_\nu\) 的定义与命题 \ref{prop:cdim-defect-measure-fourier-laplace-holographic} 的 \(t=0\) 口径，
\[
\widehat H_\nu(\xi)
=
4\pi\sum_{\nu'=1}^{r}
e^{-(1+\delta^{(\nu')})|\xi|}
\sum_{j:\delta_j=\delta^{(\nu')}}
m_j\frac{\delta^{(\nu')}}{1+\delta^{(\nu')}}e^{-\mathrm{i}\gamma_j\xi}.
\]
当 \(s=1\) 时，这正是定理 \ref{thm:cdim-comoving-horizon-scan-first-layer-extraction} 的情形。设前 \(s-1\) 层已被剥离，则由构造有
\[
\widehat H^{(s)}(\xi)
=
4\pi\sum_{\nu'=s}^{r}
e^{-(1+\delta^{(\nu')})|\xi|}
\sum_{j:\delta_j=\delta^{(\nu')}}
m_j\frac{\delta^{(\nu')}}{1+\delta^{(\nu')}}e^{-\mathrm{i}\gamma_j\xi},
\]
这与首层提取定理完全同型，只需把最小深度替换为 \(\delta^{(s)}\)。于是
\[
e^{(1+\delta^{(s)})|\xi|}\widehat H^{(s)}(\xi)
=
G_s(\xi)+\text{更快衰减项},
\]
从而第 \(s\) 层可唯一恢复。对 \(s=1,\dots,r\) 归纳，即得整个有限原子测度 \(\nu\) 的唯一恢复。证毕。
\end{proof}
因此，二维 Fourier--Laplace 单射性不仅是静态存在性命题，而且给出一个按深度递归剥离的可执行恢复程序。

\subsection{低秩缺陷的有限采样可恢复性：Hankel--Prony 双向分解}\label{subsec:cdim-hankel-prony-reconstruction}
在命题 \ref{prop:cdim-defect-measure-fourier-laplace-holographic} 的设定下，若 \(\nu\) 为 \(\kappa\) 原子测度
\[
\nu=\sum_{j=1}^{\kappa}a_j\,\delta_{(\gamma_j,\delta_j)}
\qquad(\kappa\ge 1,\ a_j>0,\ (\gamma_j,\delta_j)\ \text{两两不同}),
\]
则二维指纹化为显式有限指数和
\[
F(\xi,s)=4\pi\sum_{j=1}^{\kappa}a_j\,\frac{\delta_j}{1+\delta_j}\,e^{-s\delta_j}\,e^{-\mathrm{i}\gamma_j\xi}.
\]

\begin{proposition}[有限采样的低秩可逆性与参数分离]\label{prop:cdim-atomic-defect-hankel-prony}
固定任意 \(\xi\neq 0\) 与步长 \(\Delta s>0\)，对任意起点 \(s_0>0\) 定义等距采样序列
\(
y_n:=F(\xi,s_0+n\Delta s)
\)。
则 \((y_n)\) 为 \(\kappa\) 项指数和
\[
y_n=\sum_{j=1}^{\kappa}c_j(\xi)\,\lambda_j^{n},
\qquad
\lambda_j=e^{-\Delta s\,\delta_j},
\qquad
c_j(\xi)=4\pi a_j\frac{\delta_j}{1+\delta_j}\,e^{-s_0\delta_j}e^{-\mathrm{i}\gamma_j\xi}.
\]
因此在有限缺陷/低秩口径下，\(\{\delta_j\}\) 可由 \(\{\lambda_j\}\) 的 Prony/Hankel（矩阵 pencil）指纹恢复，而 \(\{\gamma_j\}\) 可由 \(c_j(\xi)\) 的相位项 \(e^{-\mathrm{i}\gamma_j\xi}\) 恢复；并由二维 Fourier--Laplace 的单射性保证整体可辨识。
\end{proposition}
\begin{proof}
见结论 \ref{con:xi-terminal-atomic-defect-hankel-prony}。证毕。
\end{proof}

\begin{theorem}[单频 \texorpdfstring{$2\kappa$}{2kappa} 样本的精确深度恢复]\label{thm:cdim-atomic-defect-prony-2kappa-recovery}
在本小节设定下，固定任意 \(\xi\neq 0\)、步长 \(\Delta s>0\) 与起点 \(s_0\in\RR\)。若所有深度 \(\delta_1,\dots,\delta_\kappa\) 两两不同，并记
\[
y_n:=F(\xi,s_0+n\Delta s)\qquad(n\ge 0),
\]
则
\[
y_n=\sum_{j=1}^{\kappa} c_j(\xi)\lambda_j^n,
\qquad
\lambda_j:=e^{-\Delta s\,\delta_j},
\qquad
c_j(\xi):=4\pi a_j\frac{\delta_j}{1+\delta_j}e^{-s_0\delta_j}e^{-\mathrm{i}\gamma_j\xi}.
\]
在上述深度互异条件下，仅由前 \(2\kappa\) 个样本
\[
y_0,y_1,\dots,y_{2\kappa-1}
\]
即可唯一恢复全部深度 \(\delta_j\) 与复振幅 \(c_j(\xi)\)。
\leanverified{paper\_cdim\_atomic\_defect\_prony\_2kappa\_recovery}
\end{theorem}
\begin{proof}
由上式，\((y_n)\) 是 \(\kappa\) 项指数和。定义首一湮灭多项式
\[
Q(z)=\prod_{j=1}^{\kappa}(z-\lambda_j)
=
z^\kappa+q_{\kappa-1}z^{\kappa-1}+\cdots+q_0.
\]
则 \((y_n)\) 满足长度 \(\kappa\) 的线性递推
\[
y_{n+\kappa}+q_{\kappa-1}y_{n+\kappa-1}+\cdots+q_0y_n=0
\qquad(n\ge 0).
\]
取 \(n=0,\dots,\kappa-1\)，得到关于 \(q_0,\dots,q_{\kappa-1}\) 的 \(\kappa\) 个线性方程，其系数矩阵为 Hankel 矩阵
\[
\mathcal H_\kappa=(y_{i+j})_{0\le i,j\le \kappa-1}.
\]
由于 \(\lambda_j\) 两两不同且 \(c_j(\xi)\neq 0\)，有标准 Vandermonde 分解
\[
\mathcal H_\kappa=
V\,\diag(c_1(\xi),\dots,c_\kappa(\xi))\,V^{\mathsf T},
\qquad
V_{ij}=\lambda_j^i,
\]
从而 \(\mathcal H_\kappa\) 可逆。故湮灭多项式 \(Q\) 的系数唯一，进而其根 \(\lambda_j\) 唯一确定。于是
\[
\delta_j=-\frac{1}{\Delta s}\log \lambda_j
\]
唯一恢复全部深度。
\par
最后，把 \(\lambda_j\) 代回前 \(\kappa\) 个样本方程
\[
y_n=\sum_{j=1}^{\kappa} c_j(\xi)\lambda_j^n,
\qquad
0\le n\le \kappa-1,
\]
得到 Vandermonde 线性系统。因 \(\lambda_j\) 两两不同，故其解唯一，从而全部复振幅 \(c_j(\xi)\) 也被唯一恢复。证毕。
\end{proof}
这把低秩可逆性推进为一个精确样本复杂度结论：在深度互异的 generic 情形下，单频 \(2\kappa\) 样本恰足以恢复全部深度谱。

\begin{theorem}[双频 \texorpdfstring{$4\kappa$}{4kappa} 样本的二维参数精确恢复]\label{thm:cdim-atomic-defect-two-frequency-4kappa-recovery}
在定理 \ref{thm:cdim-atomic-defect-prony-2kappa-recovery} 的假设下，再取另一频率 \(\xi'\neq \xi\)，并假设所有横向参数满足先验窗口约束
\[
\gamma_j\in I\subset\RR,
\qquad
|I|<\frac{2\pi}{|\xi-\xi'|}.
\]
则仅由两组各 \(2\kappa\) 个样本
\[
\{F(\xi,s_0+n\Delta s)\}_{n=0}^{2\kappa-1},
\qquad
\{F(\xi',s_0+n\Delta s)\}_{n=0}^{2\kappa-1},
\]
即可唯一恢复整个有限原子族
\[
\{(\gamma_j,\delta_j,a_j)\}_{j=1}^{\kappa}.
\]
若进一步处在零点重数口径 \(a_j=m_j\in\ZZ_{\ge 1}\) 下，则全部 \((\gamma_j,\delta_j,m_j)\) 亦唯一恢复。
\end{theorem}
\begin{proof}
由定理 \ref{thm:cdim-atomic-defect-prony-2kappa-recovery}，第一组样本唯一恢复 \(\{\delta_j,c_j(\xi)\}\)，第二组样本唯一恢复 \(\{\delta_j,c_j(\xi')\}\)。由于深度两两不同，可据 \(\delta_j\) 对两组结果一一配对。对每个 \(j\)，有
\[
\frac{c_j(\xi)}{c_j(\xi')}
=
e^{-\mathrm{i}\gamma_j(\xi-\xi')}.
\]
于是
\[
\gamma_j
\equiv
-\frac{1}{\xi-\xi'}\Arg\!\left(\frac{c_j(\xi)}{c_j(\xi')}\right)
\pmod{\frac{2\pi}{|\xi-\xi'|}}.
\]
由先验窗口约束 \(|I|<2\pi/|\xi-\xi'|\)，模周期歧义被唯一解除，从而 \(\gamma_j\) 唯一确定。
\par
再由
\[
|c_j(\xi)|=4\pi a_j\frac{\delta_j}{1+\delta_j}e^{-s_0\delta_j},
\]
得
\[
a_j=
\frac{|c_j(\xi)|}{4\pi}\cdot \frac{1+\delta_j}{\delta_j}\,e^{s_0\delta_j},
\]
故全部权重 \(a_j\) 唯一恢复。若 \(a_j=m_j\in\ZZ_{\ge 1}\) 表示零点重数，则同式唯一确定相应重数。证毕。
\end{proof}
因此，在最自然的 generic 条件下，二维离线零点参数的恢复不再依赖无限采样，而可以压缩到两频 \(4\kappa\) 个样本的精确口径。

\subsection{扫描剖面的有限秩核化：从边界标量到 \texorpdfstring{$\kappa$}{kappa} 维再现核 Hilbert 空间}\label{subsec:cdim-scan-profile-finite-rank-rkhs}
对有限缺陷剖面 \(H\)（合并同参原子后可写为 \(\kappa\) 项洛伦兹混合）
\[
H(x)=\sum_{k=1}^{\kappa}\frac{w_k}{(x-\gamma_k)^2+(1+\delta_k)^2},
\qquad
w_k>0,\ \delta_k>0,
\]
引入复骨架点
\(
p_k:=\gamma_k-\mathrm{i}(1+\delta_k)\in\CC_{-}
\)
与特征映射
\(
v:\RR\to\CC^\kappa
\),
\(
v(x)_k:=\frac{\sqrt{w_k}}{x-p_k}
\)。

\begin{theorem}[有限秩核的对角核化与 \texorpdfstring{$\kappa$}{kappa} 维 RKHS]\label{thm:cdim-basepoint-scan-finite-rank-rkhs}
定义核
\[
K(x,y):=\langle v(x),v(y)\rangle_{\CC^\kappa}
=\sum_{k=1}^{\kappa}\frac{w_k}{(x-p_k)(y-\overline{p_k})}
=\sum_{k=1}^{\kappa}\frac{w_k}{(x-\gamma_k+\mathrm{i}(1+\delta_k))(y-\gamma_k-\mathrm{i}(1+\delta_k))}.
\]
则 \(K\) 为正定核并满足对角恒等式
\[
\boxed{\;K(x,x)=H(x)\qquad(x\in\RR).\;}
\]
从而扫描剖面可被精确核化为一个有限秩核的对角；并且
\(\mathcal{H}:=\overline{\Span}\{v(x):x\in\RR\}\subseteq\CC^\kappa\)
在自然内积下为 \(\kappa\)-维 Hilbert 空间（等价地：\(K\) 为一维参数生成的 \(\kappa\)-维 RKHS 的再现核）。
\end{theorem}
\begin{proof}
该定理即定理 \ref{thm:xi-basepoint-scan-finite-rank-rkhs} 在本章口径下的直接译文。证毕。
\end{proof}
该核化把“边界标量剖面”提升为“有限秩核几何”，从而把层析反演的稳定性与数值审计接口严格对齐到有限秩核在采样点上的 Gram 数据恢复问题；这与 Hankel--Prony/矩阵 pencil/ESPRIT 的标准反演口径兼容。

\subsection{位复杂度屏障与共动前缀收益：二次压缩的制度成本与可审计回收}\label{subsec:cdim-bit-budget-barrier-comoving-prefix-gain}
\subsubsection{端点二次压缩与高度窗上的最坏情形位预算}
固定 Cayley 图表下，离临界线偏移在端点层余量 \(h(\rho)=1-|w_\rho|^2\) 中发生二次压缩（定理 \ref{thm:xi-offcritical-quadratic-radial-compression}）。
在高度窗 \(|\gamma|\le T\) 且偏移下界 \(\delta\ge\delta_0\) 的最坏情形中，端点余量满足显式下界（定理 \ref{thm:xi-offcritical-endpoint-resolution-lower-bound}）
\[
1-|w_\rho|^2\ \ge\ \frac{4\delta_0}{T^2+(1+\delta_0)^2},
\]
从而若以 \(2^{-p}\) 作为端点余量的最小可分辨尺度，则必有位预算下界
\[
p\ \ge\ \log_2\!\Bigl(\frac{T^2+(1+\delta_0)^2}{4\delta_0}\Bigr)
=2\log_2T+\log_2\!\frac{1}{\delta_0}+O(1)\qquad(T\to\infty),
\]
并可在精度势口径下等价写成 \(\mathcal{P}(T)=\log\pi+\log(1+T^2)\) 的显式形态（式 \eqref{eq:xi-offcritical-endpoint-bit-precision-potential-form}）。

\subsubsection{共动前缀的常数预算与制度成本}
若允许在对象层外置可寻址前缀 \(x_0\) 并取共动局部化 \(x_0=\gamma\)，则
\[
h_{\gamma}(\rho)=1-|w_{\rho,\gamma}|^2=\frac{4\delta}{(1+\delta)^2}\ \ge\ \frac{4\delta_0}{(1+\delta_0)^2}\qquad(\delta\ge\delta_0),
\]
从而端点余量阈值可在高度无关的常数位预算下证书化（推论 \ref{cor:xi-comoving-prefix-bit-budget-gain}）。
需要强调的是：该“回收”并不消失，而是被制度性转移为地址前缀预算；若前缀未给定，则相应评价在对象层不定义并读出为 \(\mathsf{NULL}\)（注 \ref{rem:xi-comoving-horizon-localization-address-null}）。

\subsubsection{紧化读出的径向位复杂度下界}
相位紧化把动态范围强制压入末位精度（式 \eqref{eq:app-precision-amplification}），并在“仅允许输出 dyadic 括号”的制度中导出硬下界：
若仅通过 \(u=\upsilon(x)=\pi^{-1}\arctan(x)\) 的 \(n\)-比特 dyadic 近似来证书化 \(x\) 且要求离线复核得到绝对误差 \(\le\varepsilon\) 的单值可核验读出，则必有（定理 \ref{thm:xi-radial-information-projection-lower-bound}）
\[
\boxed{\;
n\ \ge\ \log_2\!\Bigl(\frac{1}{\varepsilon}\Bigr)+\log_2(1+x^2)+\log_2\pi+O(1)\; }.
\]
将其应用到完成化接口的径向参数 \(x=\log|u_L(s)|=(2\Re(s)-1)\log L\) 可见：当 \(\Re(s)\neq\tfrac12\) 且允许 \(L\to\infty\) 时，证书位数不可能对 \(L\) 一致有界，必须随 \(\log_2(1+x^2)\) 增长（同上定理）。

\subsection{正交切片定理：可见模长分量、盘内零点分量与端点载荷分量}\label{subsec:cdim-orthogonal-slice-theorem}
在 unitary slice 可见域约束（\(X_{\mathrm{vis}}=\mathsf{Phase}\)）与 \(\mathbf{PW}\) 闭包类型检查下，偏离临界切片的信息在可审计意义下被强制分解为三类互不替代的正交通道：
\begin{enumerate}
  \item \textbf{协议正交（径向分量）：}任何本质上需要 \(|u|\neq 1\) 的连续径向自由度才能表达的对象，在纯相位接口中不可类型化；若拒绝外置唯一径向寄存器 \(\varrho=\log|u|\)，则读出按协议退化为 \(\mathsf{NULL}\)（定理 \ref{thm:xi-pw-no-continuous-hair}、推论 \ref{cor:xi-unique-continuous-transverse-register}，并见 \S\ref{subsec:cdim-offline-zero-claim-interface}）。
  \item \textbf{因子分解正交（盘内零点分量）：}在 Hardy/Smirnov 因子分解口径下，边界模长仅决定外因子，而盘内零点与奇异支撑完全由内因子承载；因此“离切片点状事件”的零点数据属于内因子分量，对边界模长审计结构性不可见（附录 \S\ref{subsec:unit-disk-inner-outer}）。
  \item \textbf{端点载荷正交（端点分量）：}端点邻域的伸缩半群作用分解为有限离散残差与连续精度势负载两条互不耦合的账本轴（\S\ref{subsec:cdim-endpoint-orthogonal-load}），从而端点载荷在证书闭包中对应独立记录通道，不能由提升半径采样或地址深度替代（并见 \S\ref{subsec:cdim-orthogonal-threshold-nonsubstitutable}）。
\end{enumerate}
因此，“离开 \(\Re(s)=\tfrac12\)”若要在闭环内保持可核验，只能沿上述正交不变量与外置预算轴出现；拒绝外置其中任一支撑轴，将在类型层必然触发 \(\mathsf{NULL}\) 或其可审计失败见证（定理 \ref{thm:xi-null-complete-trichotomy-offline}）。

\endinput

\subsection{判定性闭合：\texorpdfstring{$\mathsf{NULL}$}{NULL} 与其补集均具有限证书}\label{subsec:cdim-null-decidability-closure}
在 unitary slice 可见域锁定与 \(\mathbf{PW}\) 闭包类型安全下，\(\mathsf{NULL}\) 不再是“未定义/不可审计”的同义词，而被提升为一种携带\emph{有限否证证据}的结构性输出：读出为 \(\mathsf{NULL}\) 时必须能给出其失败来源的可离线复核见证；读出非 \(\mathsf{NULL}\) 时必须伴随落在可见域内的可核对地址/证书对象。

\begin{corollary}[unitary-slice 协议下的判定性闭合：\texorpdfstring{$\mathsf{NULL}$}{NULL} 的补集亦具有限证书]\label{cor:cdim-unitary-slice-decidable}
在定理 \ref{thm:xi-null-orthogonal-trichotomy-resource-nonsubstitutable} 的设定下，对任意被编译到 unitary-slice 审计接口的查询/对象 \(q\)，若其全局读出为 \(\mathsf{NULL}\)，则存在一份有限失败见证 \(\mathrm{Wit}(q)\)，其类型必落入三类正交方向之一并可由确定性验证器离线复核；若读出不为 \(\mathsf{NULL}\)，则按读出约定必伴随一份落在可见域内的可核对地址/证书对象。
因此，“可寻址性谓词”（读出是否为 \(\mathsf{NULL}\)）在该闭包制度下同时具有正向与负向的有限证书体系，从而在该制度意义下为可判定对象。
\leanverified{paper\_cdim\_unitary\_slice\_decidable}
\end{corollary}
\begin{proof}
该推论即推论 \ref{cor:xi-unitary-slice-decidable} 在本章符号下的直接译文。证毕。
\end{proof}

\subsection{正交三分律的见证类型：域外拒绝、\v{C}ech 粘合障碍与硬阈值资源不足}\label{subsec:cdim-null-witness-tags}
为与 \S\ref{subsec:cdim-null-three-way} 的三轴地址语义对齐，可将 \(\mathsf{NULL}\) 的三类正交来源抽象为失败见证的类型标签。

\begin{definition}[失败见证的类型标签]\label{def:cdim-null-witness-tag}
在 unitary slice 可见域锁定与 \(\mathbf{PW}\) 闭包语义下，称失败见证 \(\mathrm{Wit}(q)\) 的类型
\(\mathrm{tag}(\mathrm{Wit}(q))\)
取于集合
\[
\{\mathrm{Out},\ \mathrm{Cech},\ \mathrm{Hard}\},
\]
其中：
\begin{enumerate}
  \item \textbf{闭包/域外拒绝见证 \(\mathrm{Out}\)}：证明 \(q\) 的解释需要超出可见域 \(X_{\mathrm{vis}}=\mathsf{Phase}\) 的自由度（典型地需要 \(|u|\neq 1\) 的径向模式），或指称在闭包内不可建立，从而读出按协议退化为 \(\mathsf{NULL}\)（定义 \ref{def:xi-visible-read-null}；定理 \ref{thm:xi-pw-no-continuous-hair}）。
  \item \textbf{\v{C}ech 二阶粘合障碍见证 \(\mathrm{Cech}\)}：给出并核验前缀站点上的 Čech \(2\)-余圈障碍类 \([\alpha]\neq 0\)，由此推出地址纤维为空、读出为 \(\mathsf{NULL}\)（定义 \ref{def:prefix-site-h2-gluing-obstruction}；命题 \ref{prop:null-as-h2-obstruction}）。
  \item \textbf{硬阈值资源不足见证 \(\mathrm{Hard}\)}：给出 residue/端点等支撑轴上的阈值型预算违背证书，从而排除任何使读出“单值可核验”的合法外置；典型子类包括 prime-register 的分位门槛与端点二次分辨率门槛（定理 \ref{thm:xi-null-orthogonal-trichotomy-resource-nonsubstitutable} 第 (3) 类）。
\end{enumerate}
\end{definition}

\begin{proposition}[不可替代性]\label{prop:cdim-null-witness-nonsubstitutable}
上述三类失败见证在体系内相互正交、不可替代：增加某一轴的预算不能替代修复另一轴的缺陷；特别地 prime-register 的模积增长不能抵消端点分辨率门槛，任何预算提升亦不能消去 \([\alpha]\neq 0\) 的粘合障碍。
\leanverified{paper\_cdim\_null\_witness\_nonsubstitutable}
\end{proposition}
\begin{proof}
见定理 \ref{thm:xi-null-orthogonal-trichotomy-resource-nonsubstitutable}。证毕。
\end{proof}

\subsection{前缀站点与局部截面：\texorpdfstring{$\mathsf{NULL}$}{NULL} 的内禀刻画与 \v{C}ech 障碍}\label{subsec:cdim-prefix-site-local-sections-null}
在“地址先于取值”的语义下，\(\mathsf{NULL}\) 的对象论刻画是局部截面不存在。
本章已在 \S\ref{subsec:cdim-three-axis-address} 以站点与（预）层表述读出可定义性，并在命题 \ref{prop:cdim-nonnull-section-criterion} 给出严格判据：
\[
\mathrm{read}(a)\neq \mathsf{NULL}\ \Longleftrightarrow\ F(a)\neq\varnothing,
\]
并给出有限覆盖上的粘接唯一性。
在该语言下，\(\mathrm{Cech}\)-型见证对应“局部逐片存在但二阶一致性失败”，其失败由 \(H^2\) 障碍类 \([\alpha]\neq 0\) 的可核验代表给出（命题 \ref{prop:null-as-h2-obstruction}）。
\begin{theorem}[\v{C}ech--cyclotomic 量子化：\texorpdfstring{$\mathsf{NULL}$}{NULL} 的谱指纹与阶数对应]\label{thm:cdim-cech-cyclotomic-quantization}
\leanverified{paper\_cdim\_cech\_cyclotomic\_quantization}
设 \(A\) 为有限阿贝尔群，\(\widehat A:=\Hom(A,\TT)\) 为其角色群。
在固定前缀站点覆盖 \(\{b_i\}\) 上，令 \([\alpha]\in \check H^2(\{b_i\},A)\) 为定义 \ref{def:prefix-site-h2-gluing-obstruction} 的二阶粘合障碍类。
并假设存在一个由 \(A\)-标记 cocycle 诱导的有限态扩张核，使得对每个 \(\chi\in\widehat A\) 都有相应的 \(\chi\)-扭曲转移矩阵 \(B_\chi\)，其非扭曲 Perron 根为 \(\lambda:=\rho(B_{\ind})>1\)。
定义扭曲 \(\zeta\)-因子
\[
Z_\chi(t):=\det(I-tB_\chi)^{-1}.
\]
令 \(\chi_*([\alpha])\in \check H^2(\{b_i\},\TT)\) 为系数推前，并记其阶为 \(r:=\ord(\chi_*([\alpha]))\ge 1\)（即最小正整数 \(r\) 使 \(r\chi_*([\alpha])=0\)）。
则成立：
\begin{enumerate}
  \item 若 \(r>1\)，则归一化行列式
  \[
  \det\!\Bigl(I-\frac{t}{\lambda}B_\chi\Bigr)
  \]
  在单位圆上具有一个 \(r\) 阶本原单位根零点；等价地，\(\det(I-tB_\chi)\) 的零点集合在模长 \(|t|=\lambda^{-1}\) 的圆周上含有一个 \(r\) 阶本原单位根轨道。
  \item 反之，若某个非平凡 \(\chi\neq \ind\) 的 \(\det(I-\frac{t}{\lambda}B_\chi)\) 在单位圆上具有一个 \(r>1\) 阶本原单位根零点，则 \(\ord(\chi_*([\alpha]))=r\)，从而 \([\alpha]\neq 0\)，并由命题 \ref{prop:null-as-h2-obstruction} 推出该前缀地址处纤维为空，读出为 \(\mathsf{NULL}\)。
\end{enumerate}
\end{theorem}

\begin{proof}
由定义 \ref{def:prefix-site-h2-gluing-obstruction} 与命题 \ref{prop:null-as-h2-obstruction}，\([\alpha]\) 精确刻画局部截面的二阶一致性失败：\([\alpha]=0\) 当且仅当可通过重新选取局部规范使三重交上一致，从而存在全局截面；\([\alpha]\neq 0\) 则全局截面不存在并触发 \(\mathsf{NULL}\)。
\par
对任意角色 \(\chi\)，将 \(A\)-过渡数据推前到 \(\TT\) 系数后得到一个 \(\TT\)-值的胶合系统，其障碍即为 \(\chi_*([\alpha])\)。
当 \(r=\ord(\chi_*([\alpha]))>1\) 时，沿任一代表该障碍的二链回路累积相位为某个 \(r\) 阶本原单位根；该相位不能被局部规范消去。
在扭曲转移矩阵 \(B_\chi\) 的周期轨道展开中，上述相位作为乘性权重进入每条闭轨的贡献，从而在归一化谱上强制出现一个单位根相位的本征模态。
因此 \(\det(I-\frac{t}{\lambda}B_\chi)\) 在 \(|t|=1\) 上出现根，且其相位为 \(r\) 阶本原单位根，从而得到 (1)。
\par
反向地，若 \(\det(I-\frac{t}{\lambda}B_\chi)\) 在单位圆上出现 \(r>1\) 阶本原单位根零点，则存在模长为 \(\lambda\) 且相位为该单位根的特征值（或等价的谱模态）。
该模态对应一个不可被局部规范消去的 \(\TT\)-值胶合缺陷，回推得到 \(\chi_*([\alpha])\neq 0\)，且其阶与该单位根阶一致，即 \(\ord(\chi_*([\alpha]))=r\)。
于是 \([\alpha]\neq 0\)，并由命题 \ref{prop:null-as-h2-obstruction} 得到 \(\mathsf{NULL}\)。
证毕。
\end{proof}

\endinput

\subsection{从 \v{C}ech 障碍到量子结构：射影化、相位与干涉}\label{subsec:cdim-cech-projective-phase-interference}
上一小节已经把“局部可读而全局受阻”的失败机制压缩为前缀站点上的 \v{C}ech 二阶障碍类。
若仅把这一点理解为“全局截面不存在”，它仍只是一个消极陈述。
更积极地看，局部 lift 在可见层上可以投影到同一输出事件，而其粘接残差却不能被同一规范同时抹平；于是，可见统计便不再能够还原为单一路径或单一局部代表元的逐项求和。
相位、射影化与干涉因而可被理解为同一局部--全局失配机制在表示层上的不同外显，而不必被先验当作额外物理公理。
本小节沿命题 \ref{prop:cdim-nonnull-section-criterion}、命题 \ref{prop:null-as-h2-obstruction} 与定理 \ref{thm:cdim-cech-cyclotomic-quantization} 的站点语义，抽取出这一读取的最小形式。

\subsubsection{相位是粘接残差的可见投影}
在当前前缀站点口径下，可把 \(s_i\in F(b_i)\) 视作覆盖片 \(b_i\) 上的一条局部 lift（亦即局部证书对象）。
若在非空重叠 \(b_i\cap b_j\) 上，\(s_i\) 与 \(s_j\) 只在某个局部对称群 \(A\) 的作用下可比，则记过渡元为 \(g_{ij}\in A\)。
再取一个可见化通道
\[
\pi_{\mathrm{vis}}:A\longrightarrow A_{\mathrm{vis}},
\]
其中 \(A_{\mathrm{vis}}\) 为交换相位群，典型地可取为 \(\TT\) 或其有限 cyclotomic 商。
定义可见过渡元
\[
\theta_{ij}:=\pi_{\mathrm{vis}}(g_{ij})\in A_{\mathrm{vis}},
\]
以及三重交上的可见残差
\[
\omega_{ijk}:=\theta_{ij}\theta_{jk}\theta_{ki}\in A_{\mathrm{vis}}.
\]

\begin{definition}[可见过渡残差类]\label{def:cdim-visible-transition-residual-class}
沿上述记号，由局部 lift 族 \(\{s_i\}\) 所诱导的 \v{C}ech \(2\)-类
\[
[\omega]\in \check H^2(\{b_i\},A_{\mathrm{vis}})
\]
称为该 lift 族的\textbf{可见过渡残差类}。
若存在局部规范 \(h_i\in A_{\mathrm{vis}}\)，使经重定义
\[
\theta'_{ij}:=h_i\,\theta_{ij}\,h_j^{-1}
\]
之后的三重交残差在共同细化覆盖上同时平凡，则称该可见残差类在可见层上可消去。
\end{definition}

\begin{proposition}[可见相位的可消去性与 \texorpdfstring{$H^2$}{H2} 平凡性]\label{prop:cdim-visible-phase-residual-triviality}
\leanverified{paper\_cdim\_visible\_phase\_residual\_triviality}
在定义 \ref{def:cdim-visible-transition-residual-class} 的设定下，\([\omega]\) 正是定义 \ref{def:prefix-site-h2-gluing-obstruction} 中障碍类 \([\alpha]\) 经可见化通道 \(\pi_{\mathrm{vis}}\) 推前后的像。
特别地，可见过渡残差能够在同一组局部规范下被同时消去，当且仅当该推前类在 \(\check H^2(\{b_i\},A_{\mathrm{vis}})\) 中平凡。
\end{proposition}
\begin{proof}
由构造，\(\omega_{ijk}\) 正是 \(\alpha_{ijk}\) 经 \(\pi_{\mathrm{vis}}\) 推前后的可见像，因此 \([\omega]=(\pi_{\mathrm{vis}})_*[\alpha]\)。
若存在局部规范使可见过渡数据在共同细化覆盖上同步化，则新的三重交残差必为平凡余圈，故 \([\omega]=0\)。
反过来，若 \([\omega]=0\)，则可见残差仅与零余圈相差一个 coboundary；在与命题 \ref{prop:cdim-nonnull-section-criterion} 相同的共同细化口径下，可通过重新选取局部规范把可见过渡系统同步化。
证毕。
\end{proof}

\begin{corollary}[可见态的射影化]\label{cor:cdim-visible-state-projectivization}
\leanverified{paper\_cdim\_visible\_state\_projectivization}
若可见读出只依赖于局部对象在规范变换下的不变量，则可见层真正稳定的比较对象不是单个局部 lift，而是其在局部规范下的等价类。
特别地，当 \([\omega]\neq 0\) 时，绝对 lift 的选取不再具有全局可比性，留下来的仅是相对相位所确定的射影类。
这与定理 \ref{thm:cdim-cech-cyclotomic-quantization} 中由 \v{C}ech 障碍诱导的 cyclotomic 相位模态完全一致。
\end{corollary}
\begin{proof}
规范变换不改变可见输出，却改变局部代表元。
因此一切可见可区分量都必须先对局部规范取商；当 \([\omega]\neq 0\) 时，该取商后的相对相位残差不可再被抹平，从而可见对象天然只以射影类出现。
证毕。
\end{proof}

\subsubsection{同一可见端点上的多重 lift 与干涉修正}
设 \(\ell_1,\ell_2\) 为两条局部 lift，并满足
\[
\pi(\ell_1)=\pi(\ell_2)=y,
\]
其中 \(y\) 是同一可见事件。
若 \(\omega^{(1)}\) 与 \(\omega^{(2)}\) 分别为它们的可见过渡残差，则定义相对残差类
\[
[\omega(\ell_1,\ell_2)]
:=
\bigl[\omega^{(1)}(\omega^{(2)})^{-1}\bigr]
\in \check H^2(\{b_i\},A_{\mathrm{vis}}).
\]

\begin{definition}[可容许的可见统计]\label{def:cdim-admissible-visible-statistics}
称一个可见统计读法 \(W\) 是\textbf{可容许的}，若它满足：
\begin{enumerate}
  \item \textbf{局部规范不变性：}改变局部代表元但不改变其规范等价类时，\(W\) 不变；
  \item \textbf{粗粒化相容性：}先对同一可见事件作局部细分，再在可见层汇总，与直接在粗层赋值一致；
  \item \textbf{局部拼接一致性：}在有限覆盖上，若局部读法在交叠处相容，则汇总值只依赖于其全局拼接结果，而不依赖于具体细分方式。
\end{enumerate}
\end{definition}

\begin{proposition}[多 lift 干涉的最小结构]\label{prop:cdim-multilift-interference-correction}
\leanverified{paper\_cdim\_multilift\_interference\_correction}
在定义 \ref{def:cdim-admissible-visible-statistics} 的设定下，对同一可见事件 \(y\) 的两条局部 lift \(\ell_1,\ell_2\)，任一可容许统计都可写成
\[
W_y(\ell_1,\ell_2)
=
W_y(\ell_1)+W_y(\ell_2)+I_y(\ell_1,\ell_2),
\]
其中交叉项 \(I_y(\ell_1,\ell_2)\) 只依赖于 \(\ell_1,\ell_2\) 的相对可见残差类。
并且，若 \([\omega(\ell_1,\ell_2)]\neq 0\)，则 \(I_y(\ell_1,\ell_2)\) 不可能在全部规范代表元上恒等为零。
\end{proposition}
\begin{proof}
前式可取
\[
I_y(\ell_1,\ell_2)
:=
W_y(\ell_1,\ell_2)-W_y(\ell_1)-W_y(\ell_2)
\]
而得，其定义显然只涉及同一可见事件上的双 lift 数据。
由局部规范不变性与局部拼接一致性，\(W_y(\ell_1,\ell_2)\) 只能依赖于两条 lift 的相对可见残差类，而不能依赖于各自的绝对代表元；故 \(I_y\) 亦只能依赖于该相对类。

若进一步假设当 \([\omega(\ell_1,\ell_2)]\neq 0\) 时仍有 \(I_y(\ell_1,\ell_2)\equiv 0\)，则 \(W_y\) 在同一可见事件上的赋值将退化为单 lift 权重的朴素逐项相加。
但一旦相对残差类非平凡，两条 lift 之间便不存在能把它们同步到同一全局规范的共同选择；于是不同的局部细分会给出不同的逐项分解，而这些分解在可见层上对应同一个事件 \(y\)。
这与粗粒化相容性和局部拼接一致性矛盾。
故在非平凡相对残差类存在时，交叉项一般不可避免。
证毕。
\end{proof}

\subsubsection{可见态类的射影组织与后端表示接口}
\begin{definition}[可见态类]\label{def:cdim-visible-state-class}
对任意局部 lift \(\ell\)，把所有与 \(\ell\) 在局部规范变换下等价、且投影到同一可见对象的 lift 归为同一等价类，记为
\[
[\ell].
\]
称 \([\ell]\) 为一个\textbf{可见态类}。
\end{definition}

\begin{proposition}[后端线性化必须先下降到射影类]\label{prop:cdim-projective-state-organization}
\leanverified{paper\_cdim\_projective\_state\_organization}
若一个后端表示或线性化接口要与可见统计相容，则它必须先因子化经由定义 \ref{def:cdim-visible-state-class} 的商映射
\[
\ell\longmapsto [\ell]
\]
而后才能对代表元作任何线性组织。
换言之，在规范不变的可见层中，首先出现的不是绝对 lift 空间，而是其射影化后的态类空间。
\end{proposition}
\begin{proof}
由推论 \ref{cor:cdim-visible-state-projectivization}，一切可见可区分量都只能依赖于规范等价类。
因此任何仍然依赖绝对 lift 代表元的表示都不会在可见层下降为良定义对象。
于是，凡与可见统计相容的后端线性化都必须先对局部规范取商，即先经由 \([\ell]\) 而组织。
证毕。
\end{proof}

\begin{remark}[射影组织并不预设 Hilbert 化]\label{rem:cdim-projective-state-no-prior-hilbert}
命题 \ref{prop:cdim-projective-state-organization} 只说明：若后端表示存在，则它必须先下降到射影态类。
这并不先验要求该表示已经是 Hilbert 空间，也不要求具体的内积或正交分解先行给出。
在本章当前结构中，更直接可调用的后端接口仍是有限秩核化与再现核语言：见定理 \ref{thm:cdim-basepoint-scan-finite-rank-rkhs}。
\end{remark}

\subsubsection{平方型读出与 Hilbert 闭包的接口}
前述讨论已经说明两点：其一，相位首先是粘接残差的可见投影；其二，可见态首先应按射影类而不是绝对 lift 来组织。
剩余的关键问题是：何时可见统计能够进一步压缩为某种平方型读出？
在当前稿件中，这一点不再被保留为外部目标，而被后续 Hilbert 闭包小节重写为内部表示后果。

\begin{proposition}[平方型读出的 Hilbert 闭包接口]\label{conj:cdim-quadratic-readout-consistency}
\leanverified{paper\_cdim\_quadratic\_readout\_consistency}
设 \(P\) 为定义在可见态类 \([\ell]\) 上的统计读法，并要求它满足：
\begin{enumerate}
  \item \textbf{正性与归一化：}对完备粗粒化族，\(P\) 给出非负且归一的权重；
  \item \textbf{局部规范不变性：}\(P\) 不依赖于态类代表元的局部规范选择；
  \item \textbf{粗粒化相容性：}细分后再汇总与直接粗粒化给出同一统计；
  \item \textbf{张量复合封闭性：}独立系统复合后，\(P\) 与乘法结构相容；
  \item \textbf{条件期望一致性：}对局部可观测子系统的边缘化与条件期望口径一致。
\end{enumerate}
则存在某个振幅表示空间 \(V\)、一个把可见态类送入 \(V\) 的表示映射
\[
\mathcal A:[\ell]\longmapsto \mathcal A([\ell])\in V,
\]
以及 \(V\) 上的正半定二次型（在复线性实现中可取为正半定厄米型） \(Q\)，使得
\[
P([\ell])=Q\bigl(\mathcal A([\ell]),\mathcal A([\ell])\bigr)
\]
在适当归一化下成立。
更进一步，在 \S\ref{subsec:cdim-hilbert-quantum-closure} 的可读时间字、mixed-collision 正定核与附录 \ref{app:op-algebra} 的投影接口齐备时，上述二次型可被具体落实为 Hilbert 载体上的 Born 型读出 \(\Tr(\rho P)\)。
\end{proposition}
\begin{proof}
\S\ref{subsec:cdim-hilbert-quantum-closure} 的定理 \ref{thm:cdim-hilbert-carrier}、\ref{thm:cdim-born-pairing} 与 \ref{thm:cdim-hilbert-quantum-main-closure} 已把可见态类的统计读法压实为 Hilbert 载体上的投影概率配对。
因此，平方型读出在本章中已经不再是外加目标，而是后端 Hilbert 闭包的表示结论。
证毕。
\end{proof}

\begin{remark}[与现有核化接口的兼容性]\label{rem:cdim-quadratic-readout-kernel-interface}
本章已有三类后端结构与命题 \ref{conj:cdim-quadratic-readout-consistency} 相容。
其一是定理 \ref{thm:cdim-basepoint-scan-finite-rank-rkhs} 给出的有限秩核化：它把边界标量剖面重写为有限维再现核空间中的对角核读出。
其二是命题 \ref{prop:cdim-poisson-kl-all-moment-quadratic-form} 与命题 \ref{prop:cdim-poisson-kl-quadratic-form-spectral} 给出的正二次型与谱表示：它们说明在粗粒化差异与矩差通道中，平方型组织已作为内部结构出现。
其三是 \S\ref{subsec:cdim-hilbert-quantum-closure} 中由 mixed-collision 正定核、离散幺正演化与附录算子代数共同推出的 Hilbert 载体：它把上述二次型组织具体落实为量子事件与 Born 配对。
\end{remark}

\noindent
因此，在当前章的语境中，“相位--射影态--干涉--平方型读出”并非四个彼此独立的附加物理设定，而是同一 \v{C}ech 粘接失配在可见表示层上的连续闭包：先出现的是可见残差与射影态类，继而是交叉修正，最后才是由正定核与投影接口压实的 Hilbert 读出。
\noindent
这也解释了为何后续地址接口不能停留在单一可见前缀上：一旦局部 lift 只在射影意义下可比较，而交叉修正又由残差类控制，完整的可核验接口就必须把可见轴、残差轴与模式轴同时显式化，并在后端继续闭合到 \S\ref{subsec:cdim-hilbert-quantum-closure} 的量子结构。

% --- Distillation writeback ---
\begin{proposition}[可见残差的共同细化不变性]\label{prop:distill-cech-cohomology-visible-residual-common-refinement-invariance}
\textbf{Dependency status: conditional rigidity.}
沿用定义 \ref{def:cdim-visible-transition-residual-class} 的记号，并设可见系数群 \(A_{\mathrm{vis}}\) 为交换群。令
\[
\mathcal U=\{b_i\}_{i\in I}
\]
为当前有限覆盖，令
\[
\mathcal V=\{c_\alpha\}_{\alpha\in J}
\]
为其共同细化，并固定细化索引映射 \(\rho:J\to I\)，满足
\[
c_\alpha\subset b_{\rho(\alpha)}.
\]
设两族局部 lift \(\ell_1,\ell_2\) 在 \(\mathcal U\) 上分别给出可见过渡元
\[
\theta^{(r)}_{ij}\in A_{\mathrm{vis}},\qquad r=1,2,
\]
以及可见二余圈
\[
\omega^{(r)}_{ijk}:=\theta^{(r)}_{ij}\theta^{(r)}_{jk}\theta^{(r)}_{ki}\in Z^2(\mathcal U;A_{\mathrm{vis}}).
\]
定义相对可见残差
\[
\omega^{(12)}_{ijk}:=\omega^{(1)}_{ijk}\bigl(\omega^{(2)}_{ijk}\bigr)^{-1}.
\]
则：
\begin{enumerate}
  \item \(\omega^{(12)}\in Z^2(\mathcal U;A_{\mathrm{vis}})\)，且其细化拉回
  \[
  (\rho^\ast\omega^{(12)})_{\alpha\beta\gamma}
  :=\omega^{(12)}_{\rho(\alpha)\rho(\beta)\rho(\gamma)}
  \]
  给出 \(Z^2(\mathcal V;A_{\mathrm{vis}})\) 中的二余圈；
  \item 对任意两组局部规范
  \[
  h_i^{(r)}\in A_{\mathrm{vis}},\qquad r=1,2,
  \]
  所诱导的重定义
  \[
  \theta^{(r)\prime}_{ij}:=h_i^{(r)}\theta^{(r)}_{ij}\bigl(h_j^{(r)}\bigr)^{-1}
  \]
  不改变 \([\rho^\ast\omega^{(12)}]\in \check H^2(\mathcal V;A_{\mathrm{vis}})\)；
  \item 若
  \[
  [\rho^\ast\omega^{(12)}]\ne 0
  \qquad\text{in}\qquad
  \check H^2(\mathcal V;A_{\mathrm{vis}}),
  \]
  则在该共同细化覆盖上，\(\ell_1\) 与 \(\ell_2\) 的相对可见残差不能由局部规范同时消去。
\end{enumerate}
\end{proposition}
\begin{proof}
因为 \(A_{\mathrm{vis}}\) 交换，二余圈集合在逐点乘法下成群；故两个二余圈的商
\(
\omega^{(1)}(\omega^{(2)})^{-1}
\)
仍为二余圈。这给出第一条中 \(\omega^{(12)}\in Z^2(\mathcal U;A_{\mathrm{vis}})\)。细化拉回与 \v{C}ech 微分交换，即
\[
\delta(\rho^\ast\omega^{(12)})=\rho^\ast(\delta\omega^{(12)})=0,
\]
故 \(\rho^\ast\omega^{(12)}\) 是 \(\mathcal V\) 上的二余圈。

对第二条，计算单个 lift 的规范重定义。在三重交上有
\[
\omega^{(r)\prime}_{ijk}
=h_i^{(r)}\theta^{(r)}_{ij}\bigl(h_j^{(r)}\bigr)^{-1}
 h_j^{(r)}\theta^{(r)}_{jk}\bigl(h_k^{(r)}\bigr)^{-1}
 h_k^{(r)}\theta^{(r)}_{ki}\bigl(h_i^{(r)}\bigr)^{-1}
=\omega^{(r)}_{ijk}.
\]
因此 \(\omega^{(12)}\) 本身不变，其细化拉回与上同调类当然不变。

若第三条的结论失败，则存在共同细化上的局部规范，使相对残差在 \(\mathcal V\) 上被消去。按定义，这意味着 \(\rho^\ast\omega^{(12)}\) 与平凡二余圈相差一个二上边界，因而
\([
\rho^\ast\omega^{(12)}
]=0
\)
于 \(\check H^2(\mathcal V;A_{\mathrm{vis}})\)，与假设矛盾。证毕。
\end{proof}

% --- End distillation writeback ---

% --- Distillation writeback ---
\begin{definition}[有限屏幕探针塔]
\label{def:distill-grothendieck-finite-screen-probe-tower}
\textbf{依赖状态：construction.}
称一组数据
\[
\mathfrak P=
\bigl(
\mathcal L_n,Y_n,\pi_n,p_{m,n},r_{m,n}
\bigr)_{m\ge n\ge0}
\]
为一个有限屏幕探针塔，若 \(\mathcal L_n\) 与 \(Y_n\) 均为有限集合，
\[
\pi_n:\mathcal L_n\to Y_n,
\qquad
p_{m,n}:\mathcal L_m\to\mathcal L_n,
\qquad
r_{m,n}:Y_m\to Y_n
\]
满足
\[
p_{n,n}=\mathrm{id},\qquad r_{n,n}=\mathrm{id},
\]
\[
p_{m,n}\circ p_{k,m}=p_{k,n},\qquad
r_{m,n}\circ r_{k,m}=r_{k,n}
\qquad(k\ge m\ge n),
\]
以及屏幕相容性
\[
r_{m,n}\circ\pi_m=\pi_n\circ p_{m,n}.
\]
对一个兼容可见词
\[
y=(y_n)_{n\ge0}\in\varprojlim_n Y_n
\]
定义第 \(n\) 层屏幕纤维与极限完成纤维为
\[
F_n(y):=\pi_n^{-1}(y_n),
\qquad
F_\infty(y):=\varprojlim_n F_n(y).
\]
若 \(F_\infty(y)\ne\varnothing\)，则称：
\begin{enumerate}
  \item \(y\) 在 \(\mathfrak P\) 下为有限层可表对象，若 \(|F_\infty(y)|=1\) 且存在 \(n\) 使 \(|F_n(y)|=1\)；
  \item \(y\) 为严格 pro-可表对象，若 \(|F_\infty(y)|=1\) 且对所有 \(n\) 都有 \(|F_n(y)|>1\)；
  \item \(y\) 含 phantom 完成核，若 \(|F_\infty(y)|\ge2\)。
\end{enumerate}
并定义第 \(n\) 层 phantom 核为
\[
K^{\rm ph}_n(y):=
\left\{(\xi,\xi')\in F_\infty(y)^2:
\xi\ne\xi',\ \xi_n\ne\xi'_n
\right\},
\]
其中 \(\xi_n\) 表示 \(\xi\) 在 \(F_n(y)\) 中的第 \(n\) 坐标。
\end{definition}

\begin{proposition}[有限屏幕探针的可表--pro-可表--phantom 三分]
\label{prop:distill-grothendieck-finite-screen-probe-trichotomy}
\textbf{依赖状态：conditional audit.}
设 \(\mathfrak P\) 为定义 \ref{def:distill-grothendieck-finite-screen-probe-tower} 的有限屏幕探针塔，并固定
\[
y\in\varprojlim_n Y_n
\]
且 \(F_\infty(y)\ne\varnothing\)。则以下三种情形恰有一种成立：
\begin{enumerate}
  \item \(y\) 为有限层可表对象；
  \item \(y\) 为严格 pro-可表对象；
  \item \(y\) 含 phantom 完成核。
\end{enumerate}
并且第三种情形等价于存在某个有限深度 \(n\) 使
\[
K^{\rm ph}_n(y)\ne\varnothing.
\]
在定义 \ref{def:cdim-visible-state-class} 的语境中，这说明：同一可见态类的有限屏幕一致性只给出探针层对象；若 \(K^{\rm ph}_n(y)\ne\varnothing\)，则该屏幕存在有限深度不可消去的 phantom 核，不能被误读为唯一 interior lift。
\end{proposition}
\begin{proof}
由于 \(F_\infty(y)\ne\varnothing\)，其基数或者为 \(1\)，或者至少为 \(2\)。若 \(|F_\infty(y)|\ge2\)，按定义正是第三种情形。若 \(|F_\infty(y)|=1\)，则再分两种互斥情形：存在某个 \(n\) 使 \(|F_n(y)|=1\)，或对所有 \(n\) 都有 \(|F_n(y)|>1\)。前者是有限层可表对象，后者是严格 pro-可表对象。因此三种情形覆盖所有可能，且彼此互斥。

现在证明 phantom 核判据。若 \(K^{\rm ph}_n(y)\ne\varnothing\)，则存在 \(\xi\ne\xi'\in F_\infty(y)\)，故 \(|F_\infty(y)|\ge2\)，于是 \(y\) 含 phantom 完成核。反过来，若 \(|F_\infty(y)|\ge2\)，取两点 \(\xi\ne\xi'\)。作为逆极限中两条不同兼容序列，它们必在某个有限坐标 \(n\) 上不同，否则全部坐标相同便会给出 \(\xi=\xi'\)。于是 \((\xi,\xi')\in K^{\rm ph}_n(y)\)，从而某个有限深度的 phantom 核非空。证毕。
\end{proof}

% --- End distillation writeback ---

\endinput

\subsection{从射影态类到 Hilbert 型量子结构的闭包}\label{subsec:cdim-hilbert-quantum-closure}
上一小节只固定了可见相位、射影态类与干涉修正的最小局部--全局来源。
若停在这一步，量子侧仍只是一组表示约束，而尚未闭合成可直接计算的态空间、事件算子与时间演化。
本小节把该缺口压回到项目 PDF 已有的两类内部接口：其一是扫描--投影读出在稳定类型空间上的可见分布与 mixed-collision 正定核（推论 \ref{cor:pom-mixed-collision-pd-kernel}、定理 \ref{thm:conclusion-mixed-collision-characteristic-hilbert-geometry}）；其二是附录 \ref{app:op-algebra} 与 \S\ref{subsec:op-algebra-fold-quantum-channel} 已给出的投影、态与完全正可见实现。
相应的量子闭包宜分成两层理解：第一层止于抽象的、离散时间的、Hilbert 型量子力学；第二层则在此基础上，经连续时间生成元、交换可见代数的谱表示与审计局域 Euclidean 正常形，止于连续时间的局域 Euclidean 标量 Schr\"odinger 扇区。

\subsubsection{可读时间字、稳定可见分布与正定重叠核}
以下固定一个稳定可读分辨率 \(m\ge 1\)，并记相应稳定类型空间为 \(X_m\)。

\begin{definition}[可读时间字]\label{def:cdim-readable-time-word}
称有限有序列
\[
w=((a_0,t_0),\dots,(a_{r-1},t_{r-1})),
\qquad
a_j\in\mathcal A,\quad t_0<\cdots<t_{r-1}\in\mathbb Z,
\]
为一个\textbf{可读时间字}，若它满足扫描--投影生成器、稳定语法与折叠规范形的可读性要求，从而能够通过有限窗口扫描被送入稳定类型空间 \(X_m\) 并给出稳定读出。
记所有此类时间字的集合为 \(\mathcal W_m\)。
\end{definition}

\begin{proposition}[可读时间字导出稳定可见分布]\label{prop:cdim-readable-time-word-visible-distribution}
\leanverified{paper\_cdim\_readable\_time\_word\_visible\_distribution}
每个可读时间字 \(w\in\mathcal W_m\) 都导出一个稳定可见分布
\[
\mu_w\in\Delta(X_m).
\]
\end{proposition}
\begin{proof}
扫描--投影读出模型在有限窗口上把可读字送入稳定类型空间，并以折叠规范形给出统计对象。
因此每个 \(w\) 都对应于 \(X_m\) 上的一个稳定可见分布 \(\mu_w\)。
证毕。
\end{proof}

\begin{definition}[重叠核]\label{def:cdim-readable-time-word-overlap-kernel}
对任意 \(u,v\in\mathcal W_m\)，定义其重叠核为
\[
K_m(u,v):=\sum_{x\in X_m}\mu_u(x)\mu_v(x).
\]
\end{definition}

\leanverified{paper\_cdim\_readable\_time\_word\_overlap\_pd}
\begin{theorem}[稳定可见分布上的正定重叠核]\label{thm:cdim-readable-time-word-overlap-pd}
定义 \ref{def:cdim-readable-time-word-overlap-kernel} 给出的 \(K_m\) 是稳定可见分布的 pullback 正定核。
更精确地，存在特征映射
\[
\Phi_m:\Delta(X_m)\longrightarrow \ell^2(X_m),
\qquad
\Phi_m(\mu):=\bigl(\mu(x)\bigr)_{x\in X_m},
\]
使得
\[
K_m(u,v)=\langle \Phi_m(\mu_u),\Phi_m(\mu_v)\rangle_{\ell^2(X_m)}.
\]
因此，若两条可读时间字在所有重叠核配对下不可区分，则它们诱导同一稳定可见分布。
\end{theorem}
\begin{proof}
这正是推论 \ref{cor:pom-mixed-collision-pd-kernel} 在 \(q=1\) 情形下的直接应用；定理 \ref{thm:conclusion-mixed-collision-characteristic-hilbert-geometry} 进一步说明该核在稳定可见分布上是 characteristic 的。
证毕。
\end{proof}

\begin{definition}[Hilbert 预载体]\label{def:cdim-hilbert-precarrier}
在复线性包 \(\mathbb C[\mathcal W_m]\) 上定义 sesquilinear 形式
\[
\Bigl\langle \sum_i c_i u_i,\ \sum_j d_j v_j\Bigr\rangle_m
:=
\sum_{i,j}\overline{c_i}d_j\,K_m(u_i,v_j).
\]
记其零范数子空间为
\[
\mathcal N_m:=\{X\in\mathbb C[\mathcal W_m]:\langle X,X\rangle_m=0\},
\]
并定义商空间
\[
\mathcal H_{0,m}:=\mathbb C[\mathcal W_m]/\mathcal N_m.
\]
\end{definition}

\begin{theorem}[Hilbert 载体]\label{thm:cdim-hilbert-carrier}
商空间 \(\mathcal H_{0,m}\) 带有由 \(\langle\cdot,\cdot\rangle_m\) 诱导的真内积，其完备化
\[
\mathcal H_Q:=\overline{\mathcal H_{0,m}}
\]
是一个 Hilbert 空间。
\end{theorem}
\begin{proof}
由定理 \ref{thm:cdim-readable-time-word-overlap-pd}，\(\langle\cdot,\cdot\rangle_m\) 半正定。
商去零范数方向后得到真内积，而任何内积空间都有唯一 Hilbert 完备化，故 \(\mathcal H_Q\) 存在且由该内积唯一决定。
\leanverified{paper\_cdim\_hilbert\_carrier}
证毕。
\end{proof}

\subsubsection{离散幺正演化与量子事件}
\begin{definition}[一步时间移位]\label{def:cdim-one-step-time-shift}
对任意可读时间字
\[
w=((a_0,t_0),\dots,(a_{r-1},t_{r-1}))\in\mathcal W_m,
\]
定义其一步时间移位为
\[
\sigma w:=((a_0,t_0+1),\dots,(a_{r-1},t_{r-1}+1)).
\]
\end{definition}

\begin{proposition}[同步时间移位保持重叠核]\label{prop:cdim-time-shift-preserves-overlap}
对任意 \(u,v\in\mathcal W_m\)，有
\[
K_m(\sigma u,\sigma v)=K_m(u,v).
\]
\end{proposition}
\begin{proof}
附录 \ref{app:op-algebra} 已把时间平移写成
\(
P_t=U^{-t}PU^t
\)。
因此，对所有局部读出同时作同一步时间平移，只是对同一读出协议做共同共轭，不改变稳定可见分布及其 \(\ell^2\) 重叠。
故 \(K_m(\sigma u,\sigma v)=K_m(u,v)\)。
证毕。
\end{proof}
\leanverified{paper\_cdim\_time\_shift\_preserves\_overlap}

\begin{theorem}[离散幺正演化]\label{thm:cdim-discrete-unitary-evolution}
\leanverified{paper\_cdim\_discrete\_unitary\_evolution}
映射
\[
U_0:\mathcal H_{0,m}\longrightarrow \mathcal H_{0,m},
\qquad
U_0[w]:=[\sigma w],
\]
是等距双射，因而唯一延拓为 \(\mathcal H_Q\) 上的幺正算子
\[
U:\mathcal H_Q\to\mathcal H_Q.
\]
\end{theorem}
\begin{proof}
命题 \ref{prop:cdim-time-shift-preserves-overlap} 表明
\[
\langle U_0[w],U_0[v]\rangle_m
=
K_m(\sigma u,\sigma v)
=
K_m(u,v)
=
\langle [w],[v]\rangle_m.
\]
故 \(U_0\) 保持内积并下降到商空间。
由 \(\sigma^{-1}\) 的存在，\(U_0\) 还是双射，故在完备化上唯一延拓为幺正算子 \(U\)。
证毕。
\end{proof}

\begin{definition}[事件子空间]\label{def:cdim-quantum-event-subspace}
对任意地址 \(a\) 与时刻 \(t\in\mathbb Z\)，定义闭子空间
\[
\mathcal H_{a,t}
:=
\overline{\operatorname{span}}\{[w]\in\mathcal H_Q:\ w\text{ 在时刻 }t\text{ 含有地址 }a\}.
\]
\end{definition}

\begin{theorem}[量子事件]\label{thm:cdim-quantum-event-projections}
对每个地址--时间对 \((a,t)\)，存在唯一正交投影
\[
P_{a,t}:\mathcal H_Q\to\mathcal H_{a,t},
\]
并满足协变关系
\[
P_{a,t}=U^{-t}P_{a,0}U^t.
\]
\leanverified{paper\_cdim\_quantum\_event\_projections}
\end{theorem}
\begin{proof}
定义 \ref{def:cdim-quantum-event-subspace} 中的 \(\mathcal H_{a,t}\) 是 Hilbert 空间 \(\mathcal H_Q\) 的闭子空间，因此存在唯一正交投影 \(P_{a,t}\)。
又由一步时间移位的定义，\(\mathcal H_{a,t}=U^{-t}\mathcal H_{a,0}\)，从而正交投影满足
\(
P_{a,t}=U^{-t}P_{a,0}U^t
\)。
这与附录 \ref{app:op-algebra} 的时间协变公式一致。
证毕。
\end{proof}

\subsubsection{正常态、Born 配对与态空间几何}
\begin{definition}[正常态]\label{def:cdim-normal-states}
定义
\[
D(\mathcal H_Q):=
\{\rho\in\mathcal B_1(\mathcal H_Q):\ \rho\ge 0,\ \Tr\rho=1\},
\]
其中 \(\mathcal B_1(\mathcal H_Q)\) 表示迹类算子。
称 \(D(\mathcal H_Q)\) 的元素为 \(\mathcal H_Q\) 上的\textbf{正常态}。
\end{definition}

\begin{theorem}[Born 配对]\label{thm:cdim-born-pairing}
\leanverified{paper\_cdim\_born\_pairing}
对任意 \(\rho\in D(\mathcal H_Q)\) 与任意量子事件 \(P_{a,t}\)，定义
\[
\Pr_\rho(a,t):=\Tr(\rho P_{a,t}).
\]
则 \(\Pr_\rho(a,t)\in[0,1]\)，并且对任何互斥完备事件族给出概率分布。
\end{theorem}
\begin{proof}
由定理 \ref{thm:cdim-quantum-event-projections}，\(P_{a,t}\) 是正交投影，因此 \(0\le P_{a,t}\le I\)。
故
\[
0\le \Tr(\rho P_{a,t})\le \Tr(\rho)=1.
\]
若 \(\{P_i\}\) 两两正交且 \(\sum_i P_i=I\)，则
\[
\sum_i\Pr_\rho(i)
=
\sum_i\Tr(\rho P_i)
=
\Tr\Bigl(\rho\sum_iP_i\Bigr)
=
\Tr(\rho)=1.
\]
故得到概率分布。
这正把附录 \ref{app:op-algebra} 中的 \(\Tr(\rho P_a)\) 读出压成 \(\mathcal H_Q\) 上的 Born 配对。
证毕。
\end{proof}

\begin{definition}[纯态与混态]\label{def:cdim-pure-mixed-states}
称 \(\rho\in D(\mathcal H_Q)\) 为\textbf{纯态}，若它不能写成两个不同正常态的非平凡凸组合；否则称其为\textbf{混态}。
\end{definition}

\begin{proposition}[纯态的秩一刻画]\label{prop:cdim-pure-states-rank-one}
\leanverified{paper\_cdim\_pure\_states\_rank\_one}
\(\rho\in D(\mathcal H_Q)\) 是纯态，当且仅当存在单位向量 \(\psi\in\mathcal H_Q\) 使
\[
\rho=\lvert\psi\rangle\langle\psi\rvert.
\]
并且 \(\psi\) 与 \(e^{\mathrm i\theta}\psi\) 给出同一纯态。
\end{proposition}
\begin{proof}
这是密度算子凸几何的标准结论。
若 \(\rho=\lvert\psi\rangle\langle\psi\rvert\)，则
\[
\lvert e^{\mathrm i\theta}\psi\rangle\langle e^{\mathrm i\theta}\psi\rvert
=
\lvert\psi\rangle\langle\psi\rvert,
\]
故物理不变量是射线而非单位向量本身。
证毕。
\end{proof}

\begin{corollary}[波函数只是坐标表示]\label{cor:cdim-wavefunction-as-coordinate}
\leanverified{paper\_cdim\_wavefunction\_as\_coordinate}
一旦在某个表示扇区上选定一组正交基 \(\{\lvert x\rangle\}\)，纯态射线 \([\psi]\) 的坐标函数
\[
\psi(x):=\langle x\mid\psi\rangle
\]
便给出波函数表示。
因此，波函数只是纯态射线在所选基底下的坐标，而不是先验对象。
\end{corollary}
\begin{proof}
命题 \ref{prop:cdim-pure-states-rank-one} 已说明纯态由射线决定。
对任意正交基展开该射线即可得到坐标函数 \(\psi(x)\)。
证毕。
\end{proof}

\subsubsection{复合体系、纠缠与可见实现}
\begin{definition}[独立机制]\label{def:cdim-independent-mechanisms}
称两套可读机制独立，若其复合重叠核满足乘法分解
\[
K_{12}\bigl((u_1,u_2),(v_1,v_2)\bigr)
=
K_1(u_1,v_1)\,K_2(u_2,v_2).
\]
\end{definition}

\begin{theorem}[复合 Hilbert 载体]\label{thm:cdim-composite-hilbert-carrier}
若两套机制在定义 \ref{def:cdim-independent-mechanisms} 的意义下独立，则其复合 Hilbert 载体满足
\[
\mathcal H_{Q,12}\cong \mathcal H_{Q,1}\otimes\mathcal H_{Q,2}.
\]
\leanverified{paper\_cdim\_composite\_hilbert\_carrier}
\end{theorem}
\begin{proof}
在独立条件下，Gram 核满足乘法分解；因此其预 Hilbert 内积在简单张量上满足
\[
\langle [u_1]\otimes [u_2], [v_1]\otimes [v_2]\rangle
=
\langle [u_1],[v_1]\rangle\,\langle [u_2],[v_2]\rangle.
\]
由张量积 Hilbert 空间的泛性质，复合预载体的完备化唯一同构于
\(
\mathcal H_{Q,1}\otimes\mathcal H_{Q,2}
\)。
证毕。
\end{proof}

\begin{corollary}[纠缠]\label{cor:cdim-entanglement}
\leanverified{paper\_cdim\_entanglement}
若 \(\mathcal H_{Q,12}\) 上某纯态射线不能写成
\[
[\psi_1]\otimes[\psi_2],
\]
则它不是两个局域纯态的简单并置，因而给出纠缠态。
\end{corollary}
\begin{proof}
由定理 \ref{thm:cdim-composite-hilbert-carrier}，复合体系的纯态射线以张量积表示。
不能分解为简单张量的纯态射线，正是纠缠态的定义。
证毕。
\end{proof}

\begin{definition}[可见实现]\label{def:cdim-visible-realization}
称线性映射
\[
\mathcal E:\mathcal B_1(\mathcal H_Q)\to\mathcal B_1(\mathcal H_Q)
\]
为一个\textbf{可见实现}，若它保持正常态集合 \(D(\mathcal H_Q)\)。
\end{definition}

\begin{theorem}[量子信道]\label{thm:cdim-quantum-channels}
\leanverified{paper\_cdim\_quantum\_channels}
附录 \ref{app:op-algebra} 与 \S\ref{subsec:op-algebra-fold-quantum-channel} 已给出的可见实现，都可提升为完全正且迹保持的量子信道。
特别地，
\[
\mathcal E_U(\rho):=U\rho U^\ast
\]
是可见实现，而折叠量子信道提供了非幺正的可见实现族。
\end{theorem}
\begin{proof}
幺正共轭显然保持正性与迹一，因此给出可见实现。
\S\ref{subsec:op-algebra-fold-quantum-channel} 进一步在固定分辨率的可见事件扇区上显式构造了折叠量子信道、其 Kraus 表示、Choi 秩、最小环境维数与完全正性。
把该构造经由本小节的 Hilbert 表示读取到 \(\mathcal H_Q\) 上，便得到完全正且迹保持的可见实现。
证毕。
\end{proof}

\begin{corollary}[条件化更新]\label{cor:cdim-luders-update}
\leanverified{paper\_cdim\_luders\_update}
对投影测量 \(P_{a,t}\)，其后验更新写成
\[
\rho\longmapsto
\frac{P_{a,t}\rho P_{a,t}}{\Tr(\rho P_{a,t})},
\]
从而形成“先事件条件化、后经量子信道传播”的一般可见演化接口。
\end{corollary}
\begin{proof}
这正是量子信道语义中的投影测量特例。
分母由定理 \ref{thm:cdim-born-pairing} 给出，并在事件发生条件下严格为正。
证毕。
\end{proof}

\subsubsection{连续时间、连续谱表示与局域 Euclidean 扇区}
\begin{definition}[主相位生成元]\label{def:cdim-principal-phase-generator}
设 \(U\) 为定理 \ref{thm:cdim-discrete-unitary-evolution} 给出的离散幺正演化。
由谱定理，存在唯一投影值测度 \(E_U\) 使
\[
U=\int_{\TT} z\,dE_U(z).
\]
对任意 Borel 集 \(B\subseteq(-\pi,\pi]\)，定义
\[
\Pi(B):=E_U\bigl(\{e^{-\mathrm i\theta}:\theta\in B\}\bigr).
\]
定义主相位生成元
\[
H:=\int_{(-\pi,\pi]}\theta\,d\Pi(\theta).
\]
\end{definition}

\begin{proposition}[幺正的自伴对数]\label{prop:cdim-unitary-selfadjoint-log}
\leanverified{paper\_cdim\_unitary\_selfadjoint\_log}
定义 \ref{def:cdim-principal-phase-generator} 给出的 \(H\) 是有界自伴算子，满足
\[
\|H\|\le \pi,
\qquad
U=e^{-\mathrm iH}.
\]
\end{proposition}
\begin{proof}
主支相位函数 \(\theta\mapsto\theta\) 在 \((-\pi,\pi]\) 上有界且实值，因此函数演算给出 \(H\) 为有界自伴算子。
再由谱积分的定义立即得到 \(e^{-\mathrm iH}=U\)。
证毕。
\end{proof}

\begin{theorem}[连续时间幺正群]\label{thm:cdim-continuous-unitary-group}
\leanverified{paper\_cdim\_continuous\_unitary\_group}
对每个 \(s\in\mathbb R\)，定义
\[
U_s:=e^{-\mathrm isH}.
\]
则 \(\{U_s\}_{s\in\mathbb R}\) 构成唯一的强连续幺正群，满足
\[
U_0=I,\qquad U_{s+t}=U_sU_t,\qquad U_1=U.
\]
\end{theorem}
\begin{proof}
由命题 \ref{prop:cdim-unitary-selfadjoint-log}，\(H\) 为有界自伴算子，故函数演算给出幺正群 \(U_s=e^{-\mathrm isH}\)。
群性质由指数律直接成立，而有界性保证 \(s\mapsto U_s\) 在算子范数下连续，故强连续。
主相位分支一旦固定，满足 \(U_1=U\) 的连续幺正分支便唯一。
证毕。
\end{proof}

\begin{definition}[同刻可交换可见代数]\label{def:cdim-commutative-visible-algebra}
在附录 \ref{app:op-algebra} 的可交换读出口径中，定义同刻可见代数
\[
\mathfrak M_0:=W^\ast\bigl(\{P_{a,0}\}_{a\in\mathcal A}\bigr)\subseteq \mathcal B(\mathcal H_Q).
\]
\end{definition}

\begin{proposition}[同刻可见代数的交换性]\label{prop:cdim-commutative-visible-algebra}
\leanverified{paper\_cdim\_commutative\_visible\_algebra}
\(\mathfrak M_0\) 是交换 von Neumann 代数。
\end{proposition}
\begin{proof}
附录 \ref{app:op-algebra} 在可交换读出口径下把地址事件写成
\[
P_a=\prod_j Q_j(a_j),
\]
因此同刻地址事件投影两两可交换。
取其生成的 von Neumann 代数即得交换性。
证毕。
\end{proof}

\leanverified{paper\_circle\_dimension\_continuous\_spectrum\_representation}
\leanverified{paper\_cdim\_continuous\_spectrum\_representation}
\begin{theorem}[连续谱表示]\label{thm:cdim-continuous-spectrum-representation}
存在测度空间 \((X,\nu)\)、可测多重度函数 \(m(x)\) 与酉同构
\[
W:\mathcal H_Q\longrightarrow L^2(X,\nu;\mathbb C^{m(x)}),
\]
使得对任意 \(M\in\mathfrak M_0\)，存在有界可测函数 \(f_M\) 满足
\[
WMW^{-1}=M_{f_M}.
\]
特别地，对每个同刻地址事件投影 \(P_{a,0}\)，存在可测集合 \(X_a\subseteq X\) 使
\[
WP_{a,0}W^{-1}=M_{\mathbf 1_{X_a}}.
\]
\end{theorem}
\begin{proof}
命题 \ref{prop:cdim-commutative-visible-algebra} 已说明 \(\mathfrak M_0\) 为交换 von Neumann 代数。
由交换 von Neumann 代数的谱表示定理，\(\mathfrak M_0\) 在某个直接积分空间上按乘法算子表示。
对投影元而言，对应函数必为特征函数，故得到所示集合 \(X_a\)。
证毕。
\end{proof}

\begin{corollary}[坐标表示]\label{cor:cdim-coordinate-representation}
若在某个不变扇区上 \(m(x)=1\) 几乎处处成立，则该扇区上的 Hilbert 载体同构于
\[
L^2(X,\nu),
\]
并且所有同刻地址事件都按特征函数乘法作用。
\leanverified{paper\_cdim\_coordinate\_representation}
\end{corollary}
\begin{proof}
这是定理 \ref{thm:cdim-continuous-spectrum-representation} 在 multiplicity-one 情形下的直接退化。
证毕。
\end{proof}

\begin{definition}[审计局域 Euclidean 扇区]\label{def:cdim-local-euclidean-sector}
沿用命题 \ref{prop:real-input-40-logM-curvature-mismatch-spectrum}、命题 \ref{prop:real-input-40-unified-generator-alignment}、表 \ref{tab:real-input-40-local-bigeometry-invariants} 与表 \ref{tab:real_input_40_output_potential_cumulants_closed} 的记号，在 real-input \(40\)-state kernel 的 \(\theta=0\) 审计点定义
\[
A_\ast:=\Sigma^{-1}\nabla^2\log\mathfrak M(0),
\qquad
h_\ast:=A_\ast^{\mathsf T}A_\ast.
\]
称由 \((\mathbb R^3,h_\ast)\) 所张成的局域扇区为\textbf{审计局域 Euclidean 扇区}。
\end{definition}

\begin{proposition}[审计局域 Euclidean 扇区的正定性]\label{prop:cdim-local-euclidean-sector-positive}
审计局域 Euclidean 扇区是三维实内积空间；固定正交归一基后，可识别为
\[
(\mathbb R^3,h_\ast)\cong \mathbb R^3_{\mathrm{Euc}}.
\]
\end{proposition}
\begin{proof}
命题 \ref{prop:real-input-40-logM-curvature-mismatch-spectrum} 与表 \ref{tab:real-input-40-local-bigeometry-invariants} 已给出
\[
\mathrm{spec}(A_\ast)=(5.319833,\,-0.534795,\,-2.572421),
\]
三条特征值均非零，故 \(A_\ast\) 非退化。
因此对任意 \(v\neq 0\)，
\[
v^{\mathsf T}h_\ast v
=
\lVert A_\ast v\rVert^2
>
0,
\]
即 \(h_\ast\) 为正定二次型。
有限维实内积空间与 Euclidean 空间的标准线性代数同构遂得结论。
证毕。
\end{proof}
\leanverified{paper\_cdim\_local\_euclidean\_sector\_positive}

\begin{definition}[局域势能 germ]\label{def:cdim-local-potential-germ}
定义局域标量函数 \(V_\ast\) 的 germ，使其在原点满足
\[
V_\ast(0)=\log\mathfrak M(0),\qquad
\nabla V_\ast(0)=\nabla\log\mathfrak M(0),\qquad
\nabla^2V_\ast(0)=\nabla^2\log\mathfrak M(0).
\]
\end{definition}

\begin{proposition}[连续极限的二阶主符号]\label{prop:cdim-second-order-principal-symbol}
在定理 \ref{thm:cdim-continuous-unitary-group} 的连续时间分支上，若取表 \ref{tab:real_input_40_output_potential_cumulants_closed} 所审计的长时小参数极限，则生成元的主符号由压力 Hessian \(\Sigma=\nabla^2P(0)\) 的二次型给出；特别地，\(P''(0)>0\) 保证该主符号非退化。
\end{proposition}
\begin{proof}
表 \ref{tab:real_input_40_output_potential_cumulants_closed} 给出 \(P''(0)>0\) 及更高阶累积量的闭式审计值，因此压力函数在原点的展开首个非平凡偶次项为
\[
P(\theta)=P(0)+\langle\nabla P(0),\theta\rangle+\frac12\,\theta^{\mathsf T}\Sigma\theta+O(|\theta|^3).
\]
在连续时间的长时小参数缩放下，三阶及以上项进入更高阶修正，而主导项由该二次型控制；因此相应生成元在局域 scalar 扇区上的主符号正是 \(\frac12\,\xi^{\mathsf T}\Sigma\xi\)。
非退化性由 \(P''(0)>0\) 与局域 Fisher Hessian 的正定性保证。
证毕。
\leanverified{paper\_cdim\_second\_order\_principal\_symbol}
\end{proof}

\begin{theorem}[局域 Euclidean 标量 Schr\"odinger 扇区]\label{thm:cdim-local-schrodinger-sector}
在审计局域 Euclidean 扇区上，把定理 \ref{thm:cdim-continuous-unitary-group} 的连续时间生成元限制到 multiplicity-one 的局域 scalar 扇区，并去除已由平移群吸收的一阶漂移后，可得局域正常形
\[
H_{\mathrm{loc}}
=
-\frac12\Delta_{h_\ast}+V_\ast(x).
\]
在命题 \ref{prop:cdim-local-euclidean-sector-positive} 的 Euclidean 正交坐标下，上式化为
\[
H_{\mathrm{loc}}
=
-\frac12\Delta+V_\ast(x),
\]
从而纯态波函数满足局域 Schr\"odinger 方程
\[
\mathrm i\,\partial_t\psi(t,x)
=
\Bigl(-\frac12\Delta+V_\ast(x)\Bigr)\psi(t,x).
\]
\end{theorem}
\begin{proof}
命题 \ref{prop:cdim-second-order-principal-symbol} 已把局域主部固定为由 \(\Sigma\) 决定的非退化二阶符号；命题 \ref{prop:cdim-local-euclidean-sector-positive} 把该符号对应的局域几何识别为 Euclidean 内积；定义 \ref{def:cdim-local-potential-germ} 则固定了零阶局域标量项。
在 scalar 扇区中，自伴的局域二阶算子因此被压成
\(
-\frac12\Delta_{h_\ast}+V_\ast
\)
这一标准正常形。
再由定理 \ref{thm:cdim-continuous-unitary-group} 的 \(U_t=e^{-\mathrm itH}\) 对时间求导，即得局域 Schr\"odinger 方程。
证毕。
\end{proof}
\leanverified{paper\_cdim\_local\_schrodinger\_sector}

\subsubsection{主定理与边界}
\begin{theorem}[从项目 PDF 到 Hilbert 型量子结构的闭包]\label{thm:cdim-hilbert-quantum-main-closure}
\leanverified{paper\_cdim\_hilbert\_quantum\_main\_closure}
在不引入任何外部物理公理的前提下，项目 PDF 的下列内部接口
\[
\text{可读时间字}
\Longrightarrow
\text{稳定可见分布}
\Longrightarrow
\text{mixed-collision 正定核}
\Longrightarrow
\text{附录算子代数与完全正可见实现}
\]
共同推出如下两层量子闭包：
\begin{enumerate}
  \item 由 Hilbert 载体 \(\mathcal H_Q\)、量子事件投影 \(P_{a,t}\)、Born 配对 \(\Tr(\rho P_{a,t})\)、纯态与混态、波函数坐标、复合体系张量积与纠缠、离散幺正演化 \(U\) 以及一般量子信道所构成的抽象离散时间 Hilbert 型量子力学；
  \item 在第一层之上，由连续时间幺正群 \(U_t=e^{-\mathrm itH}\)、一般连续谱表示
  \[
  \mathcal H_Q\cong L^2(X,\nu;\mathbb C^{m(x)}),
  \]
  multiplicity-one 扇区上的坐标表示 \(L^2(X,\nu)\)，以及由内部审计种子控制的局域 Euclidean 标量 Schr\"odinger 扇区
\[
\mathrm i\,\partial_t\psi
=
\Bigl(-\frac12\Delta+V_\ast\Bigr)\psi
\]
  所构成的连续时间分支。
\end{enumerate}
\end{theorem}
\begin{proof}
第一层由定理 \ref{thm:cdim-hilbert-carrier}、定理 \ref{thm:cdim-quantum-event-projections}、定理 \ref{thm:cdim-born-pairing}、命题 \ref{prop:cdim-pure-states-rank-one}、推论 \ref{cor:cdim-wavefunction-as-coordinate}、定理 \ref{thm:cdim-composite-hilbert-carrier}、推论 \ref{cor:cdim-entanglement}、定理 \ref{thm:cdim-discrete-unitary-evolution} 与定理 \ref{thm:cdim-quantum-channels} 共同给出。
第二层由定理 \ref{thm:cdim-continuous-unitary-group}、定理 \ref{thm:cdim-continuous-spectrum-representation}、推论 \ref{cor:cdim-coordinate-representation}、命题 \ref{prop:cdim-local-euclidean-sector-positive}、命题 \ref{prop:cdim-second-order-principal-symbol} 与定理 \ref{thm:cdim-local-schrodinger-sector} 给出。
证毕。
\end{proof}

\begin{remark}[该闭包的精确边界]\label{rem:cdim-hilbert-quantum-boundary}
若以定理 \ref{thm:cdim-hilbert-quantum-main-closure} 的第一层终点为目标，则从可读时间字到量子信道的整条链已经不再含有未对象化的内部结构缺口；就此意义而言，本文量子侧已经闭合出抽象的、离散时间的、Hilbert 型量子力学。

第二层把这一离散结构继续延伸到连续时间幺正群、一般 \(L^2(X,\nu)\) 连续谱表示与审计局域 Euclidean 标量 Schr\"odinger 扇区。若进一步要求无条件的全局 Hamiltonian 标准形式、全局唯一的标准 \(L^2(\mathbb R^d)\) 粒子表象，或其具体质量、单位与实验参数的物理校准，则仍需另行补入更强的表示与校准分支；这些问题属于比本文更强的全局化目标，而不构成上述两层内部结构链本身的缺口。
\end{remark}

\endinput

\subsection{从同步成本到时空结构：因果锥、传播上界与连续极限}\label{subsec:cdim-sync-causal-holonomy-spacetime}
上一小节已经把相位、射影态与干涉压回到局部--全局粘接残差在统计表示层上的投影。
若把同一类残差改在同步、共同支撑与闭路复合层上读取，则时间、空间、因果与曲率也不必被当作外加几何舞台，而可被理解为协议内部可比性条件的几何重组。
本小节只抽取这一重组的最小结构：同步所需的最小延迟、共同支撑所需的最小代价、以及闭路 residual 的面积归一化极限。
因此，这里既不预设连续流形，也不预设 Lorentz 度量或固定光速常数；所谓传播上界与连续曲率，只作为这些协议量之间的表示后果出现。

\subsubsection{同步延迟、共同支撑代价与因果前序}
设 \(\mathcal X\) 为协议允许的局部状态类。
对 \(x,y\in\mathcal X\)，记 \(x\leadsto y\) 为一条允许的单步精化、同步或传输步。

\begin{definition}[同步延迟与共同支撑代价]\label{def:cdim-sync-delay-support-cost}
对任意 \(x,y\in\mathcal X\)，定义：
\begin{enumerate}
  \item \textbf{同步延迟}
  \[
  \tau_{\mathrm{sync}}(x,y)\in[0,\infty]
  \]
  为把 \(x\) 与 \(y\) 拉入同一可比较协议口径所需的最小在线延迟、最小精化深度或最小同步成本；
  \item \textbf{共同支撑代价}
  \[
  d_{\mathrm{space}}(x,y)\in[0,\infty]
  \]
  为把 \(x\) 与 \(y\) 拉入共同支撑、共同 forcing 区域或共同可比较局部图所需的最小资源代价。
\end{enumerate}
若不存在任何允许程序使该比较得以建立，则对应量记为 \(+\infty\)。
\end{definition}

\begin{definition}[因果可达前序]\label{def:cdim-causal-reachability-preorder}
对 \(x,y\in\mathcal X\)，若存在有限允许链
\[
x=x_0\leadsto x_1\leadsto \cdots \leadsto x_n=y
\]
使得
\[
\sum_{k=0}^{n-1}\tau_{\mathrm{sync}}(x_k,x_{k+1})<\infty,
\qquad
\sum_{k=0}^{n-1}d_{\mathrm{space}}(x_k,x_{k+1})<\infty,
\]
则记
\[
x\preceq y,
\]
称 \(y\) 落在 \(x\) 的\textbf{因果可达域}中。
\end{definition}

\begin{proposition}[同步与共同支撑诱导的因果前序]\label{prop:cdim-causal-preorder-from-subadditivity}
若同步延迟与共同支撑代价沿允许复合满足
\[
\tau_{\mathrm{sync}}(x,z)
\le
\tau_{\mathrm{sync}}(x,y)+\tau_{\mathrm{sync}}(y,z),
\]
\[
d_{\mathrm{space}}(x,z)
\le
d_{\mathrm{space}}(x,y)+d_{\mathrm{space}}(y,z),
\]
且在对角线上均为 \(0\)，则定义 \ref{def:cdim-causal-reachability-preorder} 所给出的关系 \(\preceq\) 构成 \(\mathcal X\) 上的一个前序。
\leanverified{paper\_cdim\_causal\_preorder\_from\_subadditivity}
\end{proposition}
\begin{proof}
反身性由零长度链给出。
若 \(x\preceq y\) 且 \(y\preceq z\)，则分别存在从 \(x\) 到 \(y\) 与从 \(y\) 到 \(z\) 的有限允许链。
将两条链首尾拼接，再利用同步延迟与共同支撑代价的次可加性，即得一条从 \(x\) 到 \(z\) 的有限允许链，其总同步成本与总支撑成本仍为有限，故 \(x\preceq z\)。
证毕。
\end{proof}

\begin{remark}[当前稿件中的同步接口]\label{rem:cdim-sync-delay-existing-interfaces}
定义 \ref{def:cdim-sync-delay-support-cost} 并非空泛记号，而已有若干后端实现。
在折叠同步模型中，最坏同步延迟由推论 \ref{cor:Ym-sync-delay} 精确给出为 \(D_{\mathrm{sync}}(m)=m-1\)；在因果右逆口径下，最小延迟与最小状态数由定理 \ref{thm:xi-time-min-delay-min-states} 同时取到；在固定窗口的概率同步口径下，最小精确同步成本又可由定理 \ref{thm:conclusion-window6-exact-fold-synchronization-mutual-information} 写成互信息极小值。
因此，同步量在本文中已经具备组合论、自动机与信息论三种可核验落点。
\end{remark}

\subsubsection{传播上界作为同步--支撑比的极大不变量}
若一条允许步既伴随有限同步延迟，又伴随有限共同支撑代价，则其空间--时间比给出该步的局部传播速率。

\begin{definition}[局部传播比与传播上界]\label{def:cdim-local-propagation-ratio}
对任一允许的基本传播步
\[
e:x\leadsto y
\qquad
\text{且}\qquad
0<\tau_{\mathrm{sync}}(x,y)<\infty,
\]
定义其\textbf{局部传播比}为
\[
\rho(e)
:=
\frac{d_{\mathrm{space}}(x,y)}{\tau_{\mathrm{sync}}(x,y)}.
\]
进一步定义
\[
c_*:=\sup\bigl\{\rho(e): e\ \text{为允许的基本传播步}\bigr\}\in[0,\infty].
\]
称 \(c_*\) 为该协议族的\textbf{传播上界}。
\end{definition}

\begin{proposition}[链式传播上界]\label{prop:cdim-chainwise-propagation-bound}
设
\[
x=x_0\leadsto x_1\leadsto\cdots\leadsto x_n=y
\]
为任一允许链，并定义其总同步延迟与总共同支撑代价
\[
\tau_{\mathrm{chain}}
:=
\sum_{k=0}^{n-1}\tau_{\mathrm{sync}}(x_k,x_{k+1}),
\qquad
d_{\mathrm{chain}}
:=
\sum_{k=0}^{n-1}d_{\mathrm{space}}(x_k,x_{k+1}).
\]
若定义 \ref{def:cdim-local-propagation-ratio} 中的 \(c_*\) 有限，且链上每一步满足
\[
\tau_{\mathrm{sync}}(x_k,x_{k+1})>0
\quad\text{或}\quad
d_{\mathrm{space}}(x_k,x_{k+1})=0,
\]
则恒有
\[
d_{\mathrm{chain}}\le c_*\,\tau_{\mathrm{chain}}.
\]
\leanverified{paper\_cdim\_chainwise\_propagation\_bound}
\end{proposition}
\begin{proof}
由 \(c_*\) 的定义，对每一步 \(x_k\leadsto x_{k+1}\) 都有
\[
d_{\mathrm{space}}(x_k,x_{k+1})
\le
c_*\,\tau_{\mathrm{sync}}(x_k,x_{k+1}).
\]
对 \(k\) 求和即得所示不等式。
\end{proof}

\begin{remark}[局域可构造边与全局锥边界的分层]\label{rem:cdim-local-edge-vs-global-cone-layering}
这里需要区分两种不同层级的几何对象。单条允许步 \(x\leadsto y\) 及其同步延迟、共同支撑代价，给出的是局域可构造边；它只说明协议内部存在一条合法传输或同步步骤。相反，\(c_*\)、\(J^\pm\)、\(I^\pm\) 与 \(N^\pm\) 则属于全局刚性层，因为它们要在所有允许步与所有有限链上取上确界、闭包与边界。也就是说，因果锥边界不是一根先验画好的线，而是从大量局域可构造边经比例控制与链式闭包后才显现出来的全局不变量。
\end{remark}

\begin{definition}[未来锥、类时区与边界候选]\label{def:cdim-causal-cone-candidates}
对任意 \(x\in\mathcal X\)，定义其未来与过去锥候选为
\[
J^+(x):=\{y\in\mathcal X:\ x\preceq y\},
\qquad
J^-(x):=\{y\in\mathcal X:\ y\preceq x\}.
\]
若存在一条从 \(x\) 到 \(y\) 的允许链以及某个 \(\varepsilon>0\)，使
\[
d_{\mathrm{chain}}
\le
(c_*-\varepsilon)\tau_{\mathrm{chain}},
\]
则称 \(y\) 属于 \(x\) 的\textbf{类时未来候选}，记为 \(I^+(x)\)；并定义传播边界候选
\[
N^+(x):=J^+(x)\setminus I^+(x).
\]
过去类时区与边界 \(I^-(x),N^-(x)\) 同理定义。
\end{definition}

\begin{conjecture}[传播锥的观察者协变性]\label{conj:cdim-causal-cone-observer-covariance}
设 \(T:\mathcal X\to\mathcal X\) 为一个允许的观察者变换，并满足：
\begin{enumerate}
  \item \(T\) 把允许的局部传播步送到允许的局部传播步；
  \item \(T\) 保持同步延迟与共同支撑代价在同一规范账本下的等价类；
  \item \(T\) 保持局部细化、粗粒化与张量复合的可比较性结构。
\end{enumerate}
则预期有：
\begin{enumerate}
  \item \(x\preceq y\) 当且仅当 \(Tx\preceq Ty\)；
  \item 传播上界 \(c_*\) 在 \(T\) 下不变；
  \item \(T(J^\pm(x))=J^\pm(Tx)\)、\(T(I^\pm(x))=I^\pm(Tx)\)、\(T(N^\pm(x))=N^\pm(Tx)\)。
\end{enumerate}
换言之，观察者变化首先应保持的不是某个既定坐标表达，而是因果可达性、锥内部与锥边界的分层。
\end{conjecture}

\subsubsection{小回路 holonomy、归一化离散曲率与连续极限}
若上一小节中的可见相位来自局部 lift 在重叠上的粘接残差，则几何侧的曲率应当来自同一类残差沿小回路复合后的累积。
在本文当前稿件中，这一点已经在 dyadic holonomy 的边界重建与通量缺陷理论里获得了严格后端。

\begin{definition}[dyadic 归一化 holonomy 曲率场]\label{def:cdim-dyadic-normalized-holonomy-curvature}
沿用定理 \ref{thm:conclusion-nonabelian-dyadic-stokes-energy-reconstruction} 的设定。
对边长 \(h_m=2^{-m}\) 的 dyadic 方形剖分 \(\mathcal Q_m\)，若
\[
\operatorname{Hol}_A(\partial Q)=\exp(\Omega(Q))
\qquad (Q\in\mathcal Q_m),
\]
则定义分片常值的\textbf{归一化离散曲率场}
\[
\mathscr R_m(x)
:=
\sum_{Q\in\mathcal Q_m}
\frac{\Omega(Q)}{h_m^2}\,\mathbf 1_Q(x).
\]
它把小回路 holonomy 的对数主项写成面积归一化后的局部曲率密度。
\end{definition}

\begin{corollary}[dyadic holonomy 归一化曲率的连续极限]\label{cor:cdim-dyadic-holonomy-curvature-limit}
在定理 \ref{thm:conclusion-nonabelian-dyadic-stokes-energy-reconstruction} 的设定下，有
\[
\mathscr R_m\longrightarrow F_{12}
\qquad\text{于 }L^2(R_1),
\]
并且
\[
\|\mathscr R_m\|_{L^2(R_1)}^2
\longrightarrow
\int_{R_1}|F_{12}(x)|^2\,\mathrm d x.
\]
因此，连续曲率对象在该 dyadic 后端上可被解释为归一化小回路 holonomy 的 \(L^2\)-极限。
\end{corollary}
\begin{proof}
在定理 \ref{thm:conclusion-nonabelian-dyadic-stokes-energy-reconstruction} 的证明中，\(\mathscr R_m\) 正是其中记作 \(\mathscr H_m\) 的分片常值场。
该证明已经给出
\[
\|\mathscr R_m-P_mF_{12}\|_{L^2(R_1)}\longrightarrow 0,
\]
其中 \(P_mF_{12}\) 为 \(F_{12}\) 关于 dyadic \(\sigma\)-代数的条件期望。
又 \(P_mF_{12}\to F_{12}\) 于 \(L^2(R_1)\)，故由三角不等式得
\[
\mathscr R_m\to F_{12}\qquad\text{于 }L^2(R_1).
\]
第二式即定理 \ref{thm:conclusion-nonabelian-dyadic-stokes-energy-reconstruction} 对 \(\mathcal E_m^{\mathrm{NA}}(A)=\|\mathscr R_m\|_2^2\) 的陈述。
\end{proof}

\leanverified{paper\_cdim\_dyadic\_holonomy\_curvature\_limit}

\begin{remark}[通量缺陷给出的后验证书]\label{rem:cdim-dyadic-holonomy-flux-defect-certificate}
定理 \ref{thm:conclusion-nonabelian-dyadic-flux-defect-telescoping} 进一步表明：对任意 dyadic 方形 \(Q_0\)，粗尺度 holonomy 对数 \(\Omega(Q_0)\) 与连续曲率通量 \(\int_{Q_0}F_{12}\) 之差，可以完全重写为一列绝对收敛的 dyadic 通量缺陷级数。
因此，连续曲率不仅是小回路 holonomy 的极限对象，而且其偏离量本身也可由纯边界数据驱动的多尺度证书链后验审计。
\end{remark}

\begin{conjecture}[holonomy 细化极限的普适曲率候选]\label{conj:cdim-holonomy-curvature-refinement-limit}
设 \(\{\mathcal U_h\}_{h\downarrow 0}\) 为一族细化覆盖，且每个尺度 \(h\) 上都给定局部过渡数据与闭路 holonomy，使得：
\begin{enumerate}
  \item 小回路对数 holonomy 在尺度 \(h\) 上具有一致的主支选择与局部线性化；
  \item 面积归一化后的离散曲率场在紧集上弱预紧；
  \item 离散 Bianchi 型约束与细化过程相容；
  \item 多尺度通量缺陷满足与定理 \ref{thm:conclusion-nonabelian-dyadic-flux-defect-telescoping} 同型的可求和控制。
\end{enumerate}
则应存在子列 \(h_j\downarrow 0\)，使相应的归一化离散曲率场弱收敛到某个连续曲率候选 \(R\)；并且若极限唯一，则 \(R\) 只由离散 holonomy 的细化极限决定。
这表明连续曲率不必预先给定，而应当作为局部粘接残差在小回路复合层上的尺度稳定首项出现。
\end{conjecture}

\noindent
因此，在本章当前语境下，时间、空间、因果与曲率都可以被压回到同一条协议链上：同步延迟决定时间尺度，共同支撑代价决定空间比较，二者的极大传播比决定因果锥的边界，而小回路 holonomy 的面积归一化极限则给出连续曲率候选。
这条链仍未恢复唯一的 Lorentz 型度量，也未导出场方程；但它已经说明，所谓时空骨架并不位于协议之外，而是由同步、传播、比较与闭路 residual 共同选出的几何组织。
在本章的后续编译接口中，这也意味着地址对象不能只记录单一可见前缀，而必须继续把可见轴、残差轴与模式轴一起显式化，以承载同步账本、传播约束与 residual 的可核验外置。

\endinput


\subsection{三轴前缀与地址完备性：可见轴、残差轴与模式轴}\label{subsec:cdim-three-axis-prefix-completeness}
定义 \ref{def:multi_axis_address} 将完备地址规范打包为
\(
a=(a_{\mathrm{vis}},a_{\mathrm{res}},a_{\mathrm{mode}})
\)；
其中 \(a_{\mathrm{vis}}\) 负责可见域柱前缀，\(a_{\mathrm{res}}\) 负责同像碰撞的单射化外置，而 \(a_{\mathrm{mode}}\) 承载协议允许的正交模式轴（包括径向寄存器与最小实现维数稳定轴）。
在该打包下，“读出非 \(\mathsf{NULL}\)”需要三轴同时通过过滤与阈值核验；任一轴缺失或被禁止均触发与其支撑轴对齐的失败见证（\S\ref{subsec:cdim-null-witness-tags}），并在 \S\ref{subsec:cdim-address-compiler} 的编译—核验流程中被离线复核。

\subsection{可靠性/完备性的分离与二歧性接口：扩展保存的制度化}\label{subsec:cdim-soundness-completeness-dichotomy}
在 \(\mathbf{PW}\) 闭包类型安全下，“可寻址 \(\Rightarrow\) 可审计”属于可靠性（soundness）方向；而“离切片零点事件 \(\Rightarrow\) 在体系内可寻址”属于完备性（completeness）方向，必须由额外可审计条件补齐。
对离切片零点断言而言，闭包制度给出严格二歧性：要么外置唯一径向寄存器以扩展保存，要么输出有限失败见证并读出为 \(\mathsf{NULL}\)。

\begin{corollary}[“第四寄存器二歧性”的严格化：离切片零点断言在体系内只有两种可证伪形态]\label{cor:cdim-fourth-register-dichotomy}
在 unitary slice 可见域锁定 \(X_{\mathrm{vis}}=\mathsf{Phase}\)（定义 \ref{def:xi-visible-read-null}）、\(\mathbf{PW}\) 闭包类型安全（定理 \ref{thm:xi-offset-null-type-safety}）与“无未记录连续轴”原则（定理 \ref{thm:xi-pw-no-continuous-hair}）下，任何声称存在离切片零点（需要 \(|u|\neq 1\) 的径向自由度；等价地在圆盘化变量中需要 \(|w|<1\) 的盘内信息）的断言，在协议语义中只能以如下二者之一出现：
\begin{enumerate}
  \item \textbf{显式外延模式轴：}在地址多轴前缀中加入径向记录 \(\varrho(u)=\log|u|\)，从而把域外自由度对象化为可审计输出；
  \item \textbf{有限失败见证：}输出一份离线可核验的失败见证，证明该断言在既定闭包/可见域制度下不可类型化，其读出必退化为 \(\mathsf{NULL}\)。
\end{enumerate}
并且不存在第三条路径：任何试图以“与相位解耦的额外连续坐标寄存器”绕过 \(\mathsf{NULL}\) 的方案，要么在闭包中被拒绝并读出为 \(\mathsf{NULL}\)，要么在因子意义下等价于上述唯一径向外延 \(\varrho(u)\)。
\leanverified{paper\_cdim\_fourth\_register\_dichotomy}
\end{corollary}
\begin{proof}
该推论即推论 \ref{cor:xi-fourth-register-dichotomy} 在本章语境下的直接译文；并与 \S\ref{subsec:cdim-offline-zero-claim-interface} 的二歧性接口表述一致。证毕。
\end{proof}

\subsection{完备性缺口的可审计分解：维数稳定证书与配额证书}\label{subsec:cdim-completeness-gap-audit-decomposition}
在 unitary-slice 的零点审计口径中，“读出为 \(\mathsf{NULL}\)”的体系性来源可被压缩为若干后端核验任务：模式轴侧的维数稳定证书与 residue 侧的配额证书是其中两类不可替代的核心接口。
\par
具体而言，模式轴完备性可由 Hankel 秩一致有界或等价的 ACC-Rank 稳定性刻画（定理 \ref{thm:xi-completeness-iff-hankel-bounded}；注 \ref{rem:xi-acc-rank}），而同像碰撞方向的单射化预算由 prime-register 的分位门槛强制（推论 \ref{cor:pom-prime-register-quantile-budget}；并见定理 \ref{thm:address-defect-law} 的最坏情形下界）。
两类证书与端点分辨率门槛（定理 \ref{thm:xi-offcritical-endpoint-resolution-lower-bound}）一道，构成 \S\ref{subsec:cdim-address-compiler} 中“单射化预算与障碍检查”的离线输入接口。

\subsection{\texorpdfstring{$\mathsf{NULL}$}{NULL}—守恒闭环：将空值来源编译为门条件与缺陷账本}\label{subsec:cdim-null-conservation-loop}
在本章的可见域锁定与扩展保存语义下，“信息守恒”并不等价于“永不出现 \(\mathsf{NULL}\)”；
其严格含义是：每一种 \(\mathsf{NULL}\) 来源都必须被编译为可审计的门条件与缺陷账本项，并在必要时通过外置正交记录轴恢复可逆性。
正交三分律与不可替代性（定理 \ref{thm:xi-null-orthogonal-trichotomy-resource-nonsubstitutable}；定理 \ref{thm:cdim-budget-failure-separation}）正是该闭环的类型化表达。

% --- Distillation writeback ---
% --- Distillation writeback: Wang--Zahl finite-test Frostman/slab dichotomy ---
\begin{definition}[有限测试 Frostman 骨架与 slab 归一化资料]\label{def:distill-cdim-frostman-slab-data}
\textbf{依赖状态：construction.}
固定分辨率 \(m\)。令 \(\mathcal A_m\) 为已经通过本节 unitary-slice 预过滤的有限地址族。令 \(Q_m\) 为有限投影词测试族；每个 \(q\in Q_m\) 是映射
\[
q:\mathcal A_m\longrightarrow \Pi_q
\]
到某个有限集合 \(\Pi_q\)。对非空 \(E\subset\mathcal A_m\)，定义本段局部使用的投影词复杂度
\[
\dim_{Q_m}(E):=\sum_{q\in Q_m}\bigl(|q(E)|-1\bigr)\in\mathbb Z_{\ge 0}.
\]
给定有理数 \(0<\eta<1\)，称 \(E\) 满足有限测试 Frostman 非集中，记作
\[
E\in\mathsf{Sk}^{\rm frost}_m(\eta;Q_m),
\]
若对所有 \(q\in Q_m\) 与 \(y\in\Pi_q\) 均有
\[
|E\cap q^{-1}(y)|\le \eta |E|.
\]
若相反存在 \((q,y)\) 使 \(|E\cap q^{-1}(y)|>\eta |E|\)，则称
\[
E_{q,y}:=E\cap q^{-1}(y)
\]
为 \(E\) 的 \(q\)-slab。

令 \(\mathsf{Bad}^{\varepsilon}_m(E)\) 表示 \(E\) 对当前 \(\mathsf{NULL}\) 闭包审计命题构成 \(\varepsilon\)-近极值反例。称 \(Q_m\) 配备了针对 \(\mathsf{Bad}^{\varepsilon}_m\) 的 slab 归一化资料，若对每个 slab \(E_{q,y}\) 给定有限像
\[
N_{q,y}(E_{q,y})=:E'\subset\mathcal A_m
\]
以及截面 \(\iota_{q,y}:E'\to E_{q,y}\)，满足 \(N_{q,y}\circ\iota_{q,y}=\mathrm{id}_{E'}\)，并满足下列四项：
\begin{enumerate}
  \item \textbf{重入：} \(E'\) 仍属于同一个 unitary-slice 审计接口，且 \(\mathrm{Out},\mathrm{Cech},\mathrm{Hard}\) 三类失败见证在 \(E'\) 上有定义。
  \item \textbf{见证拉回：} \(E'\) 上任一有限失败见证可沿 \(\iota_{q,y}\) 拉回为原族 \(E\) 上同类型的有限失败见证；若 \(E'\) 给出非 \(\mathsf{NULL}\) 地址证书，则其与 \(\iota_{q,y}\) 复合给出 \(E\) 的局部地址证书。
  \item \textbf{复杂度单调：} 对所有 \(r\in Q_m\) 有 \(|r(E')|\le |r(E_{q,y})|\le |r(E)|\)，并且若 \(|q(E)|>1\)，则
  \[
  \dim_{Q_m}(E')\le \dim_{Q_m}(E)-1.
  \]
  \item \textbf{反例局部性：} 若 \(\mathsf{Bad}^{\varepsilon}_m(E)\) 且 \(E_{q,y}\) 为 slab，则 \(\mathsf{Bad}^{\varepsilon}_m(E')\)。
\end{enumerate}
\end{definition}

\begin{theorem}[最小近极值反例的 Frostman--slab 二分]\label{thm:distill-cdim-frostman-slab-minimal-counterexample}
\textbf{依赖状态：closure.}
沿定义 \ref{def:distill-cdim-frostman-slab-data}，并假设 slab 归一化资料存在。令 \(B\subset\mathcal A_m\) 为非空有限族，满足 \(\mathsf{Bad}^{\varepsilon}_m(B)\)，且在所有可由 \(B\) 经过有限次 slab 归一化得到并仍满足 \(\mathsf{Bad}^{\varepsilon}_m\) 的族中，\(B\) 对字典序
\[
(\dim_{Q_m}(E),|E|)
\]
最小。则必有以下二者之一：
\begin{enumerate}
  \item \(B\in\mathsf{Sk}^{\rm frost}_m(\eta;Q_m)\)；
  \item 存在 \(q\in Q_m\) 使 \(|q(B)|=1\)。此时 \(B\) 已塌缩到一个投影词 slab 终端扇区，且该扇区上的任意有限见证均可由定义 \ref{def:distill-cdim-frostman-slab-data} 的拉回机制提升为原地址族的有限见证。
\end{enumerate}
更一般地，从任意 \(\mathsf{Bad}^{\varepsilon}_m(B_0)\) 出发，反复选择非 Frostman 的 slab 并归一化；每一次非终端归一化都使 \(\dim_{Q_m}\) 严格下降。因此该过程在至多 \(\dim_{Q_m}(B_0)\) 次非终端步骤后，到达有限测试 Frostman 骨架或某个投影词 slab 终端扇区。
\end{theorem}
\begin{proof}
若 \(B\in\mathsf{Sk}^{\rm frost}_m(\eta;Q_m)\)，则为第一种情形。否则，按定义存在 \(q\in Q_m\) 与 \(y\in\Pi_q\)，使
\[
B_{q,y}=B\cap q^{-1}(y),\qquad |B_{q,y}|>\eta |B|.
\]
令 \(B':=N_{q,y}(B_{q,y})\)。由反例局部性，\(\mathsf{Bad}^{\varepsilon}_m(B')\) 仍成立；由重入条件，\(B'\) 与 \(B\) 属于同一审计接口，所以可与 \(B\) 在相同的最小性序中比较。

若 \(|q(B)|>1\)，复杂度单调给出
\[
\dim_{Q_m}(B')\le \dim_{Q_m}(B)-1,
\]
这与 \(B\) 在可达反例族中的最小性矛盾。因此在非 Frostman 情形下只能有 \(|q(B)|=1\)，即第二种终端 slab 情形。

对一般初始族 \(B_0\)，只要当前族既非有限测试 Frostman 又非终端 slab，就可选取一个 slab 并归一化；上述复杂度单调使非负整数 \(\dim_{Q_m}\) 严格下降。故非终端步骤数至多为初始复杂度。若过程中产生失败见证或非 \(\mathsf{NULL}\) 地址证书，见证拉回条件保证它可逐层提升到原族；若没有提前停止，则有限下降必在 Frostman 骨架或终端 slab 处停止。证毕。
\end{proof}

% --- End distillation writeback ---

% --- Distillation writeback: slab rank certificate ---
\begin{proposition}[slab 低复杂度的投影词维数与 Hankel 秩证书]\label{prop:distill-cdim-slab-hankel-rank-certificate}
\textbf{依赖状态：rigidity.}
沿定义 \ref{def:distill-cdim-frostman-slab-data}。对 \(E\subset\mathcal A_m\) 与 \(q\in Q_m\)，记
\[
d_q(E):=|q(E)|.
\]
则 \(d_q(E)\le D\) 是有限投影词子范畴维数证书：证书只需列出不超过 \(D\) 个元素组成的集合 \(S\subset\Pi_q\)，并逐项核验 \(q(E)\subset S\)。

进一步，设该地址族 \(E\) 具有有限 Markov 表示：存在有限标记图 \((\Gamma_E,\widetilde\ell_E)\)、一个有限闭路审计族 \(\mathcal C_E\)、以及满射 \(\rho:\mathcal C_E\twoheadrightarrow E\)，使得 \(q(\rho(\gamma))\) 由闭路 \(\gamma\) 的投影词标记决定。令 \(K\subset\mathbb R\) 为子域，并给定单射 \(\beta:q(E)\to K\) 与正权 \(a_x\in K_{>0}\)（\(x\in E\)）。定义
\[
h_r:=\sum_{x\in E}a_x\beta(q(x))^r\qquad(r\ge 0),
\qquad
H_D(E,q):=(h_{i+j})_{0\le i,j\le D}.
\]
则对每个整数 \(D\ge 1\)，
\[
d_q(E)\le D
\quad\Longleftrightarrow\quad
\operatorname{rank}_K H_D(E,q)\le D.
\]
若 \(E\) 来自标记时间 \(\zeta\) 的有限系数窗口，则这些矩 \(h_0,\dots,h_{2D}\) 可由命题 \ref{prop:distill-xi-sector-marked-hankel-detector} 的标记迹系数读出。因此，投影词 slab 塌缩既可由 \(d_q(E)\) 的有限枚举证明，也可由精确 Hankel 秩下降证明。
\end{proposition}
\begin{proof}
第一项是有限集合核验：给定 \(S\)，逐个 \(x\in E\) 计算 \(q(x)\) 并检查其属于 \(S\)，即可证明 \(|q(E)|\le |S|\le D\)。反向若 \(d_q(E)\le D\)，取 \(S=q(E)\) 即得证书。

对 Hankel 判别，将 \(E\) 按 \(q\)-值分组。因 \(\beta\) 在 \(q(E)\) 上单射，\(\beta(q(E))\) 的不同取值数正好为 \(d_q(E)\)。每个取值 \(\lambda\) 的聚合权重为
\[
A_\lambda:=\sum_{x\in E,\ \beta(q(x))=\lambda}a_x,
\]
且由 \(a_x>0\) 得 \(A_\lambda>0\)。于是 \(h_r=\sum_\lambda A_\lambda\lambda^r\)，这正是命题 \ref{prop:distill-xi-sector-marked-hankel-detector} 的 Hankel 型矩序列。该命题给出 \(\operatorname{rank}_K H_D(E,q)\le D\) 当且仅当不同 \(\lambda\) 的个数至多为 \(D\)，即 \(d_q(E)\le D\)。若闭路审计族 \(\mathcal C_E\) 由标记时间 \(\zeta\) 的有限系数窗口切出，则同一命题的迹系数公式给出所需矩的有限读出。证毕。
\end{proof}

% --- End distillation writeback ---

% --- End distillation writeback ---

\endinput


\subsection{Jensen 缺陷的完全有限化：盘内零点计数的半径族证书}\label{subsec:cdim-jensen-defect-finiteization}
令 \(F_\zeta\) 为圆盘化完成化函数（附录 \ref{app:unit-circle-phase-gate}），并对 \(0<\varrho<1\) 定义 Jensen 缺陷 \(D_\zeta(\varrho)\)（定义 \ref{def:app-jensen-defect}）。
由 Jensen--Blaschke 求和式（命题 \ref{prop:app-jensen-formula-defect}），若 \(\{a_k\}\) 为 \(F_\zeta\) 在 \(|w|<\varrho\) 内的零点（按重数计），则
\[
D_\zeta(\varrho)=\sum_{|a_k|<\varrho}\log\!\left(\frac{\varrho}{|a_k|}\right)\ \ge\ 0,
\qquad
D_\zeta(\varrho)=0\ \Longleftrightarrow\ F_\zeta\ \text{在 }|w|<\varrho\text{ 内无零点}.
\]
因此，“开圆盘内无零点”可被压缩为一族半径参数的圆周平均缺陷约束，且该压缩与零点计数严格等价、不依赖数值近似假设。
与此同时，有限半径网格核验不可避免地产生结构性盲区：若仅在有限半径集合上评估缺陷，则盘内零点只有在穿越采样阈值后才进入求和项，从而存在不可判别窗口（推论 \ref{cor:app-jensen-finite-sampling-blindspot}）。

\subsection{幂覆盖的缺陷共变性与尺度交换律：重参数化与固定半径显影}\label{subsec:cdim-power-cover-scale-exchange}
\paragraph{幂半群在 Jensen 层面的严格共变性}
对圆盘内解析函数 \(F\) 及其幂拉回 \(F^{(m)}(w):=F(w^m)\)，Jensen 缺陷满足严格共变
\[
D_{F^{(m)}}(\varrho)=D_F(\varrho^m)
\qquad(0<\varrho<1,\ m\ge 1)
\]
（命题 \ref{prop:app-jensen-defect-power-covariance}）。
因此幂半群 \(w\mapsto w^m\) 在 Jensen 通道上不引入新的连续自由度，而仅对半径变量施加确定性重参数化 \(\varrho\mapsto \varrho^m\)。

\paragraph{尺度交换律：端点边界层缺陷的固定半径显影}
当离切片零点被二次压缩到端点边界层时，幂覆盖与 Adams--Norm 折叠门的复合协议可将该不可见缺陷搬运到任意预设的紧子圆盘并在固定半径处显影：对任意 \(r_\star\in(0,1)\) 存在显式尺度 \(m_\star\) 使 \(|w_\rho|^{m_\star}\le r_\star\)（式 \eqref{eq:xi-adams-norm-mstar-def}--\eqref{eq:xi-adams-norm-mstar-radius}），并在任意固定 \(R\in(r_\star,1)\) 处得到高度无关的 Jensen 正缺陷下界
\[
D_{\mathrm{Nm}_{m_\star}[F_\zeta]}(R)\ \ge\ \log\!\frac{R}{r_\star}\ >\ 0
\]
（推论 \ref{cor:xi-adams-norm-fixed-radius-jensen-certificate}）。
该交换机制把端点边界层中的径向不可见性外置为尺度参数，从而使“固定半径采样”的盲区可被制度性穿透。

\paragraph{端点通道的线性增益与正交性}
需要强调的是：幂拉回在 Jensen 缺陷通道仅作半径重参数化，而在端点通道上可出现可核验的尺度放大。
在“偏移仅由 Blaschke 数据承载”的口径下，端点吸收系数 \(\Phi_{-1}(B)\)（命题 \ref{prop:xi-endpoint-absorption-coefficient}）在奇数 Adams 门下满足线性重标定 \(\Phi_{-1}(B^{(m)})=m\Phi_{-1}(B)\)（定理 \ref{thm:xi-endpoint-absorption-adams-rescaling}），从而同一尺度参数 \(m\) 在径向平均通道上不产生增益而在端点通道上提供独立的放大机制；该对照与尺度交换律的审计口径相容（定理 \ref{thm:xi-symmetry-audit-tension-scaled}）。

\subsection{反演对偶与谱侧共轭退化：幺正切片上的自对偶闭合}\label{subsec:cdim-inversion-duality-spectral-conjugation}
在乘法 Haar 测度 \(\dd^{\times}x=\dd x/x\) 的 Hilbert 空间 \(L^{2}_{\times}:=L^{2}((0,\infty),\dd^{\times}x)\) 上，定义反演对偶算子
\[
(Jf)(x):=\overline{f(1/x)}.
\]
由于 \(\dd^{\times}x\) 在 \(x\mapsto 1/x\) 下不变，\(J\) 为反线性等距对合（\(J^{2}=I\)）；Lebesgue 口径下的等距反演仅差一个 Jacobian 归一化因子（备注 \ref{rem:mellin-inversion-jacobian}）。
与带临界移位的 Mellin 读出 \(\mathcal{M}\) 相容地，反演对偶在谱侧满足如下共轭恒等式。

\begin{proposition}[Mellin 反演对偶：谱侧共轭退化]\label{prop:cdim-mellin-inversion-duality}
\leanverified{paper\_cdim\_mellin\_inversion\_duality}
对足够良好的 \(f\) 与任意 \(s\in\CC\)，成立恒等式
\[
\boxed{\;(\mathcal{M}(Jf))(s)=\overline{(\mathcal{M}f)(1-\overline{s})}\;}
\]
从而在幺正切片 \(s=\tfrac12+\mathrm{i}t\) 上有谱侧逐点共轭退化
\[
\boxed{\;(\mathcal{M}_{1/2}(Jf))(t)=\overline{(\mathcal{M}_{1/2}f)(t)}\qquad(t\in\RR).\;}
\]
\end{proposition}
\begin{proof}
见命题 \ref{prop:mellin-conjugation-identity}。证毕。
\end{proof}

\subsection{视界熵缺陷的单标量判据：盘内零点事件的边界压缩}\label{subsec:cdim-horizon-defect-scalar-criterion}
在圆盘完成化口径下，Jensen 缺陷 \(D_{\zeta}(\varrho)\) 将“盘内零点禁入”压缩为单参标量的边界行为（定义 \ref{def:app-jensen-defect}；并见 \S\ref{subsec:cdim-jensen-defect-finiteization}）。

\begin{theorem}[RH 的单标量判据：视界纯度]\label{thm:cdim-horizon-purity-criterion}
\leanverified{paper\_cdim\_horizon\_purity\_criterion}
以下命题等价：
\begin{enumerate}
  \item[(i)] RH 成立；
  \item[(ii)] \(\displaystyle \lim_{\varrho\uparrow 1}D_{\zeta}(\varrho)=0\)；
  \item[(iii)] \(\displaystyle \limsup_{\varrho\uparrow 1}D_{\zeta}(\varrho)=0\)。
\end{enumerate}
\end{theorem}
\begin{proof}
见定理 \ref{thm:xi-horizon-purity-criterion}。证毕。
\end{proof}

\subsection{排斥半径与缺陷上界链：证书序列的趋一机制}\label{subsec:cdim-repulsion-radius-chain}
对任意 \(0<\varrho<1\) 定义排斥半径
\[
r_{\ast}(\varrho):=\varrho\exp\!\bigl(-D_{\zeta}(\varrho)\bigr)\in(0,\varrho].
\]

\begin{theorem}[缺陷诱导的排斥半径]\label{thm:cdim-defect-repulsion-radius}
对任意 \(0<\varrho<1\)，\(F_{\zeta}\) 在子盘 \(|w|<r_{\ast}(\varrho)\) 内无零点。
更定量地，若 \(D_{\zeta}(\varrho)\le\varepsilon\)，则 \(r_{\ast}(\varrho)\ge \varrho e^{-\varepsilon}\)。
并且，若存在 \(\varrho_n\uparrow 1\) 使 \(r_{\ast}(\varrho_n)\uparrow 1\)，则 \(F_{\zeta}\) 在 \(\mathbb{D}\) 内无零点，从而 RH 成立。
\end{theorem}
\begin{proof}
见定理 \ref{thm:dephys-defect-repulsion-radius} 与推论 \ref{cor:app-jensen-repulsion-radius-to-one}。证毕。
\end{proof}
\leanverified{paper\_cdim\_defect\_repulsion\_radius}

\subsection{对数半径导数律：缺陷导数即盘内零点计数}\label{subsec:cdim-defect-log-derivative}
记盘内零点计数函数（按重数计）
\[
N_{F_{\zeta}}(\varrho):=\#\{a:\ F_{\zeta}(a)=0,\ |a|<\varrho\}.
\]

\begin{proposition}[对数半径导数律与驱动律]\label{prop:cdim-jensen-defect-log-derivative}
对几乎处处的 \(\varrho\in(0,1)\)，有
\[
\boxed{\;
\frac{\dd}{\dd(\log\varrho)}D_{\zeta}(\varrho)=N_{F_{\zeta}}(\varrho)\; }.
\]
并且对几乎处处的 \(\varrho\in(0,1)\)，排斥半径满足驱动律
\[
\boxed{\;
\frac{\dd}{\dd(\log\varrho)}\log r_{\ast}(\varrho)=1-N_{F_{\zeta}}(\varrho)\; }.
\]
\leanverified{paper\_cdim\_jensen\_defect\_log\_derivative}
\end{proposition}
\begin{proof}
见命题 \ref{prop:app-jensen-defect-log-derivative} 与推论 \ref{cor:dephys-repulsion-radius-zero-count-driving}。证毕。
\end{proof}

\subsection{盘内零点的不可消去边界印记：阈值化下界与闭式半径律}\label{subsec:cdim-jensen-single-zero-imprint}
若 \(F_{\zeta}\) 在 \(\mathbb{D}\) 内存在零点 \(a\) 满足 \(|a|<1\)，则对任意 \(\varrho\in(|a|,1)\) 有
\[
D_{\zeta}(\varrho)\ge \log\!\left(\frac{\varrho}{|a|}\right),
\qquad
\liminf_{\varrho\uparrow 1}D_{\zeta}(\varrho)\ge \log(1/|a|)>0
\]
（命题 \ref{prop:app-jensen-single-zero-lower-bound}）。
特别地，对离临界线零点 \(\rho=\beta+\mathrm{i}\gamma\)（\(\beta<\tfrac12\)）所诱导的圆盘像 \(w_{\rho}\)，对任意 \(\varrho\in(|w_{\rho}|,1)\) 有
\[
D_{\zeta}(\varrho)\ge \log\!\left(\frac{\varrho}{|w_{\rho}|}\right).
\]
并且 \(|w_{\rho}|\) 具有完全闭式：由 \eqref{eq:app-offcritical-radius-closed} 可得
\[
|w_{\rho}|^{2}=\frac{\gamma^{2}+(1-\delta)^{2}}{\gamma^{2}+(1+\delta)^{2}},
\qquad
\delta:=\tfrac12-\beta>0,
\]
从而端点二次压缩律与阈值化缺陷下界可被显式对齐。

\subsection{谱侧共轭—视界缺陷—排斥半径：一条完全闭合的等价证书链}\label{subsec:cdim-spectral-conjugation-horizon-defect-repulsion-chain}
本节把幺正切片的谱侧共轭封口（\S\ref{subsec:cdim-inversion-duality-spectral-conjugation}）、视界纯度的单标量压缩（\S\ref{subsec:cdim-horizon-defect-scalar-criterion}）与排斥半径的禁入律（\S\ref{subsec:cdim-repulsion-radius-chain}）组织为一条实现无关的闭合证书链：等距读出锁定切片位置，缺陷的边界极限锁定 RH，而有限半径缺陷上界给出明确的子盘零点排除半径，并可被完全离线复核。

\subsubsection{幺正切片与谱侧共轭退化：范数守恒读出的对偶封口}
在 \(L^2_{\times}=L^{2}((0,\infty),\dd^{\times}x)\) 上，Mellin--Plancherel 幺正切片定理（定理 \ref{thm:mellin-unitary-slice-unique}）给出唯一等距读出：\(\mathcal{M}_{1/2}:L^{2}_{\times}\to L^{2}(\RR,\dd t/2\pi)\) 幺正。
并且任意偏离切片 \(\Re(s)=\sigma\) 的读出必然引入权重扭曲 \(x^{\sigma-\frac12}\)（式 \eqref{eq:mellin-sigma-weight-twist}），等价于在连续对象层外置一条额外权重记录轴。
\par
在同一语义中，反演对偶算子 \(J\) 的谱侧共轭退化由命题 \ref{prop:cdim-mellin-inversion-duality} 给出；因此，若要求读出既保持 \(L^{2}\)-范数又在对偶下闭合，则协议在调和分析层面被结构性地约束到幺正切片，而无需引入新的连续寄存器。

\subsubsection{视界纯度判据：RH 的单标量极限等价}
圆盘完成化 \(F_{\zeta}\) 的 Jensen 缺陷 \(D_{\zeta}(\varrho)\) 由定义 \ref{def:app-jensen-defect} 给出，并满足 Jensen--Blaschke 求和式（命题 \ref{prop:app-jensen-formula-defect}），从而 \(D_{\zeta}(\varrho)\ge 0\) 且
\[
D_{\zeta}(\varrho)=0\ \Longleftrightarrow\ F_{\zeta}\ \text{在 }|w|<\varrho\text{ 内无零点}.
\]
在此基础上，定理 \ref{thm:cdim-horizon-purity-criterion} 把 RH 完全压缩为单参标量函数的视界极限：
\[
\mathrm{RH}\ \Longleftrightarrow\ \lim_{\varrho\uparrow 1}D_{\zeta}(\varrho)=0.
\]
进一步地，推论 \ref{cor:app-rh-iff-jensen-defect-zero} 给出完全有限化表述：RH 等价于对一切 \(\varrho\in(0,1)\) 皆有 \(D_{\zeta}(\varrho)=0\)；而附录定理 \ref{thm:app-rh-iff-jensen-defect-vanishing-sequence} 表明只需存在一列 \(\varrho_n\uparrow 1\) 使 \(D_{\zeta}(\varrho_n)\to 0\) 即可推出 RH。

\subsubsection{排斥半径与禁入律：缺陷上界推出子盘零点排除}
由缺陷定义出排斥半径 \(r_{\ast}(\varrho)=\varrho\exp(-D_{\zeta}(\varrho))\)（\S\ref{subsec:cdim-repulsion-radius-chain}）。
定理 \ref{thm:cdim-defect-repulsion-radius} 给出严格禁入结论：在任意半径层 \(\varrho\) 上，有限数据对 \((\varrho, D_{\zeta}(\varrho))\) 确定一个显式子盘 \(|w|<r_{\ast}(\varrho)\)，其中零点事件被完全排除；并且当 \(D_{\zeta}(\varrho)\le \varepsilon\) 时有定量下界 \(r_{\ast}(\varrho)\ge \varrho e^{-\varepsilon}\)。
若存在 \(\varrho_n\uparrow 1\) 使 \(r_{\ast}(\varrho_n)\uparrow 1\)，则 \(F_{\zeta}\) 在 \(\mathbb{D}\) 内无零点，从而 RH 成立。

\subsection{缺陷—零点计数的微分等价：折角、重数与排斥半径的驱动律}\label{subsec:cdim-defect-zero-count-differential-equivalence}
Jensen 缺陷不仅提供半径族的静态禁入证书，还携带盘内零点计数的动态账本信息。
命题 \ref{prop:cdim-jensen-defect-log-derivative} 给出缺陷对 \(\log\varrho\) 的导数几乎处处等于盘内零点计数 \(N_{F_{\zeta}}(\varrho)\)，从而缺陷曲线的折角、斜率与排斥半径的增长率均可被零点半径分布完全驱动。

\subsubsection{对数半径导数律：缺陷的导数即盘内零点计数}
命题 \ref{prop:cdim-jensen-defect-log-derivative} 表明对几乎处处的 \(\varrho\in(0,1)\) 有
\[
\frac{\dd}{\dd(\log\varrho)}D_{\zeta}(\varrho)=N_{F_{\zeta}}(\varrho),
\]
并因此给出缺陷的单调性与零点进入半径层时的可核验跃迁结构（推论 \ref{cor:app-jensen-defect-kink-structure}）。

\subsubsection{折角定位与视界深度：阈值半径的不可去原子}
对 \(r\in(0,1)\)，记 \(m(r)\) 为满足 \(|a|=r\) 的盘内零点总重数，则推论 \ref{cor:app-jensen-defect-kink-structure} 给出斜率跳变律：把 \(D_{\zeta}\) 视为 \(\log\varrho\) 的函数，则在 \(\varrho=r\) 处左右导数之差等于 \(m(r)\)。
因此，每个零点在半径轴上留下一个不可去的折角原子，其位置正是阈值半径 \(r=|a|\)。
\par
对离临界线零点 \(\rho=\beta+\mathrm{i}\gamma\)（\(\beta<\tfrac12\)）的圆盘像 \(w_{\rho}\)，令 \(r_{\rho}:=|w_{\rho}|\)。
由 \eqref{eq:app-offcritical-radius-closed} 的闭式，
\[
r_{\rho}^{2}=\frac{\gamma^{2}+(1-\delta)^{2}}{\gamma^{2}+(1+\delta)^{2}},
\qquad
\delta:=\tfrac12-\beta>0,
\]
并定义视界深度
\[
h(\rho):=1-r_{\rho}^{2}=\frac{4\delta}{\gamma^{2}+(1+\delta)^{2}}.
\]
当 \(|\gamma|\to\infty\) 且 \(\delta\) 有界时，有 \(h(\rho)\sim 4\delta/\gamma^{2}\)，表明离切片信息在圆盘接口中被二次压入边界层，从而其折角阈值半径 \(r_{\rho}\) 趋近于 \(1\) 并触发有限采样盲区（\S\ref{subsec:cdim-boundary-blindspot-loop}）。

\subsubsection{排斥半径的零点计数驱动律：增长率与盘内零点的等价阻塞}
由 \(r_{\ast}(\varrho)=\varrho e^{-D_{\zeta}(\varrho)}\) 与命题 \ref{prop:cdim-jensen-defect-log-derivative} 的驱动律，几乎处处成立
\[
\frac{\dd}{\dd(\log\varrho)}\log r_{\ast}(\varrho)=1-N_{F_{\zeta}}(\varrho).
\]
从而一旦盘内出现零点（存在 \(\varrho\) 使 \(N_{F_{\zeta}}(\varrho)\ge 1\)），则在其后的半径层上 \(\log r_{\ast}\) 的增长率不再为正，排斥半径不可能随 \(\varrho\uparrow 1\) 严格趋一；该阻塞机制与定理 \ref{thm:cdim-defect-repulsion-radius} 的“趋一证书 \(\Rightarrow\) RH”结论相容。

\subsection{有限采样盲区、端点二次压缩与可审计的两类破盲路径}\label{subsec:cdim-finite-sampling-blindspot-endpoint-compression-deblind}
本节将有限半径采样的不可判别性（盲区）与端点二次压缩的标度制度并置，并指出在协议内解除盲区的两条可审计路径：共动前缀局部化与可逆层析通道。

\subsubsection{有限半径采样的结构性盲区：未越过阈值半径则离线零点不可判别}
给定有限采样半径集合 \(\{\varrho_j\}_{j=1}^{M}\subset(0,1)\)，记 \(\varrho_{\max}:=\max_{1\le j\le M}\varrho_j\)。
推论 \ref{cor:app-jensen-finite-sampling-blindspot} 表明缺陷采样族 \(\{D_{\zeta}(\varrho_j)\}\) 仅能反映满足 \(|a|<\varrho_{\max}\) 的盘内零点；任何满足 \(|a|>\varrho_{\max}\) 的盘内零点在该有限采样上严格不可见。
特别地，若存在离临界线零点 \(\rho\) 的圆盘像 \(w_{\rho}\) 满足 \(\varrho_{\max}\le |w_{\rho}|\)，则仅凭该有限采样无法排除该零点；要使其在采样上被强制显现，必须 \(\varrho_{\max}>|w_{\rho}|\)。
\par
在高高度区间（\(|\gamma|\to\infty\)、\(\delta\) 有界）下，视界深度满足 \(1-|w_{\rho}|^{2}\sim 4\delta/\gamma^{2}\)（\S\ref{subsec:cdim-defect-zero-count-differential-equivalence}）。
因此，有限采样的分辨率裕度 \(1-\varrho_{\max}\) 若不达到 \(\delta/\gamma^{2}\) 量纲，则该离线零点在采样上不可检出。

\subsubsection{闭环（压缩 \texorpdfstring{$\Rightarrow$}{=>} 盲区 \texorpdfstring{$\Rightarrow$}{=>} 破盲）：前缀共动或可逆层析}
端点二次径向压缩把离切片偏移推入边界层，从而在固定图表制度与有限半径采样下不可避免地产生盲区（\S\ref{subsec:cdim-boundary-blindspot-loop}）。
解除盲区存在两条可审计路径：其一是外置可寻址前缀并实施共动局部化，使高度红移被消解为与偏移 \(\delta\) 同阶的线性可见量（命题 \ref{prop:cdim-comoving-defect-implies-delta-bound}，推论 \ref{cor:xi-comoving-prefix-bit-budget-gain}）；其二是进入可逆层析通道，把二维参数编码为一维指数谱并可由 Fourier 反演恢复（定理 \ref{thm:cdim-comoving-horizon-scan-fourier-inversion}）。

\subsubsection{地址超度量与分辨率屏障：共享前缀深度即半径收缩}
在 dyadic 地址模型下，地址超度量 \(d_{\mathrm{addr}}\) 的闭球恰为 dyadic 盒：共享 \(b\) 位前缀等价于在深度 \(b\) 的证书层级上不可分离（\S\ref{subsec:cdim-phase-prefix-quadratic-barrier}）。
在零点密度与端点压缩的合成下，分辨率预算存在不可绕过的下界：命题 \ref{prop:cdim-phase-prefix-occupancy-barrier} 给出，若要求高度带 \((T,2T]\) 上的相位地址映射在 \(b\) 位精度下可达单射，则必须
\[
b\ \ge\ 2\log_2 T+\log_2\log T-O(1).
\]
该标度由平均零点间距 \(\asymp 2\pi/\log T\) 与相位门端点导数的 \(T^{-2}\) 压缩共同决定；因此，不外置前缀时位预算主项的 \(2\log_2T\) 不可由仅提升固定图表分辨率替代。

\endinput
